\documentclass[11pt,reqno]{amsart}
\usepackage[T1]{fontenc}
\usepackage{lmodern,amsmath,amssymb,mathtools,microtype,booktabs,array,longtable}
\usepackage[margin=1in]{geometry}
\usepackage{graphicx}
\usepackage{xr-hyper}
\usepackage[hidelinks]{hyperref}
\ifdefined\pdfinfoomitdate\pdfinfoomitdate=1\fi
\ifdefined\pdftrailerid\pdftrailerid{}\fi
\ifdefined\pdfsuppressptexinfo\pdfsuppressptexinfo=15\fi
\numberwithin{equation}{section}
\newtheorem{theorem}{Theorem}[section]
\newtheorem{proposition}[theorem]{Proposition}
\newtheorem{lemma}[theorem]{Lemma}
\newtheorem{corollary}[theorem]{Corollary}
\theoremstyle{definition}\newtheorem{definition}[theorem]{Definition}
\newtheorem{example}[theorem]{Example}
\theoremstyle{remark}\newtheorem{remark}[theorem]{Remark}
\newcommand{\E}{\mathbb E}\newcommand{\Prb}{\mathbb P}
\newcommand{\R}{\mathbb R}\newcommand{\one}{\mathbf1}
\newcommand{\Q}{\mathcal Q}\newcommand{\TV}{\mathrm{TV}}
\newcommand{\cp}{\mathrm{cp}}\newcommand{\ev}{\mathrm{ev}}\newcommand{\st}{\mathrm{st}}
\DeclareMathOperator{\tr}{tr}\DeclareMathOperator{\dist}{dist}
\DeclareMathOperator{\law}{Law}\DeclareMathOperator{\rank}{rank}
\DeclareMathOperator{\supp}{supp}
\emergencystretch=2em

\externaldocument[][nocite]{supplement-refs}[General_Theta_Foundations_I_R78_Supplement.pdf]
\hypersetup{pdftitle={General Theta Foundations I: Acquired Geometry and Causal Resource Transfer},pdfauthor={Qian Qi}}
\title[Acquired geometry and causal resource transfer]{General Theta Foundations I:\\Acquired Geometry and Causal Resource Transfer}
\author{Qian Qi}\date{October 2026}
\subjclass[2020]{Primary 60J10, 68Q45; Secondary 20C15, 52A20, 60G40, 62C20, 81P15, 90C22, 94A29}
\keywords{causal experiment, predictive state, positive realization, minimal memory, finite group, quantum instrument, convex duality}
\begin{document}
\begin{abstract}
We study finite positive memory in causal experiments and the resource
needed to implement optimal experimental policies.  For every finite
reset instrument with positive call costs and no closed failure set,
we determine the complete memory-constrained all-initial-state cost
frontier by an attained finite variational principle.  Every controller
with finite worst-initial-state cost is dominated, at every initial state and with the same command
rows, by a controller whose next label is a deterministic function of
the current reset report.  The principle uses only one scalar potential
per reset state, allows singular failure effects and arbitrary optimal
ties, and gives exact synthesis for algebraic data.  The frontier is
continuous in the primitive instrument.  An explicit rational family
has an exact cost formula at every label budget and arbitrarily many
strict resource improvements, which persist under noncommuting
perturbations.  The inherited shared-preparation theory supplies a
common permutation normal form, exact initialization charges and joint
Fourier memory laws.  Operational factorization certificates also give
unrestricted-decoder approximation lower bounds for algebraic qubit
experiments.  General causal comparison, predictive compression,
finite-order measurement, rank-constrained discrimination, variational
duality and local spectral geometry retain their separate hypotheses
and complete proofs.
\end{abstract}
\maketitle
\section{Introduction}\label{sec:current-overview}\label{sec:introduction}

A predictive state records the information in an experimental past that
can affect future laws.  Its linear dimension measures one form of
compression.  Implementing those laws with finite positive memory also
requires stochastic transitions, a common decoder and an admissible
initialization.  Implementing an optimal policy introduces a further
question: how much does the best attainable cost change when the device
is allowed one additional persistent label?  Predictive dimension,
the number of physical states and the number of optimal decision values
need not answer that question.

This paper relates predictive, convex and spectral structure to finite
positive implementations of prepared causal experiments.  A controller
has one fixed parameter-independent program.  It sees only the declared
reports, and retains no uncounted history, clock or private seed.
Persistent labels, the protected physical state, program-table bits and
working arithmetic storage are separate resources.  A prescribed pure
initial label is counted.  Freely prepared mixed starts are identified
whenever that different convention is used.

The central result is an exact finite-memory variational principle for
reset instruments.  It covers every memory budget, singular failure
effects, arbitrary finite Hilbert dimension, several commands,
state-dependent costs and failure probabilities, and tied optimal
actions.  The proof removes arbitrary stochastic memory updates without
adding a row.  An explicit family then solves every memory budget and
has arbitrarily long strict resource hierarchies.  These results are
developed alongside the general shared-preparation realization theory,
operational approximation lower bounds and the original causal
comparison.  The theorem classes and their quantifiers remain distinct.

\subsection*{The complete reset memory frontier}

Let $a\in\mathcal A$ be a command and $i\in\{1,\ldots,s\}$ a failure
report.  The physical failure branch is
\[
 \mathcal I_{a,i}(\rho)=\tr(E_{a,i}\rho)\sigma_i,
 \qquad E_{a,i}\succeq0,\quad \sum_iE_{a,i}\preccurlyeq I.
\]
The report itself prepares the known state $\sigma_i$.  The observer
does not receive an additional state measurement or free reset.  A
successful call must terminate and failure forbids termination.  The
mean call cost is $\tr(G_a\rho)$ with
$c_*I\preccurlyeq G_a\preccurlyeq c^*I$ and $c_*>0$; a realized cost is
not an extra observation.  The initial density matrix is arbitrary and
unknown.  Write $V_M$ for the least worst-initial-state expected cost
over autonomous controllers with at most $M$ live labels and a prescribed
pure initial label.

Absence of a closed failure set is necessary and sufficient for a
finite-cost controller, even with one randomized command row
(Theorem~\ref{thm:reset-control}).  Put
\[
 Q_{iaj}=\tr(E_{a,j}\sigma_i),\qquad c_{ia}=\tr(G_a\sigma_i).
\]
For a row palette $p^1,\ldots,p^M\in\Delta(\mathcal A)$ arising from
a controller with finite worst-initial-state cost, consider the finite
decision problem whose actions are exactly these rows.  Its
least Bellman potential is
\[
 x_i=\min_{1\le k\le M}
           \sum_a p_a^k\left(c_{ia}+\sum_jQ_{iaj}x_j\right).
\]
Theorem~\ref{thm:r78-palette-normal-form} proves that every finite-cost
controller is dominated, separately at every initial density matrix,
by one with the same rows, the same initial row $p^1$, and deterministic
report-only routing $i\mapsto\kappa(i)$.  The argument uses the
all-initial-state requirement at an essential point: starting the
original controller with its fixed initial label from each $\sigma_i$
gives a finite comparison policy for every state of the row-palette
problem.  It does not assert transience of unreachable pairs in the
original product of physical and controller states.

Let $C_0$ be an explicit finite worst-state cost of a full-support
one-row controller.  Theorem~\ref{thm:r78-reset-frontier} gives
\begin{equation}\label{eq:r78-intro-frontier}
\begin{split}
 V_M=\min\{b:\;&p^k\in\Delta(\mathcal A),\quad
      \kappa:\{1,\ldots,s\}\to\{1,\ldots,M\},\quad
      x\in[0,C_0]^s,\quad b\in[0,C_0],\\
 &x_i\ge\sum_a p_a^{\kappa(i)}
                  (c_{ia}+\sum_jQ_{iaj}x_j)\quad(1\le i\le s),\\
 &\sum_a p_a^1(G_a+\sum_i x_iE_{a,i})\preccurlyeq bI\}.
\end{split}
\end{equation}
The minimum is attained and each feasible point constructs an actual
controller.  Conversely, every cost bound below $C_0$ realized by any
stochastic controller gives such a point.  Only $s$ scalar continuation
potentials are needed, independently of $M$; no $sM$-state product
resolvent is imposed.  Singular effects and arbitrary Bellman ties are
allowed.  For exact algebraic primitive data, this compact finite union
of semialgebraic sets decides every cost threshold and gives an exact
minimizer.  No polynomial running time or executed general elimination
is asserted.

Proposition~\ref{prop:r78-frontier-stability} proves continuity of each
$V_M$ in the admissible primitive instrument.  The lower argument uses
the bounded potentials and compact row palettes; the upper argument
uses stability of a fixed proper realizing controller.  Thus a strict
finite resource gap persists under sufficiently small perturbations,
including perturbations that remove commutation among effects and
cost operators.

There are exact growing-dimensional instances.  For any nonempty proper
action subsets $A_1,\ldots,A_s\subset\{1,\ldots,q\}$, the inherited
rational construction has $E_{a,i}=(r/s)I$ and costs in $[1,4]$.
Theorem~\ref{thm:r78-cover-frontier} computes $V_M$ by a finite minimum
over at most $M-1$ continuing actions and a nonempty initial support.
Its optimum has pure continuing rows and an initial row uniform on
that support.  This purity conclusion is for this displayed family;
the general normal form preserves arbitrary mixed command rows.
For $A_i=\{i\}$, $s=q=n\ge2$ and $r=1/2$, the full formula becomes
\begin{equation}\label{eq:r78-intro-hierarchy}
 V_M=4+\frac1n+
      \frac{(n-1)\max\{n-M+1,0\}}{n^2},\qquad M\ge1.
\end{equation}
It gives $V_1>V_2>\cdots>V_{n+1}=V_*$ in dimension $2n$, with
one exact decrement $(n-1)/n^2$ at every step.  In this family the full
gap is exactly twice the older first-residual certificate $D_M$.
The distinction between a lower certificate and the true constrained
optimum is therefore quantitative, not only terminological.

\subsection*{The optimum threshold and positive exposure}

The unrestricted reset optimum is the spectral minimum
$V_*=\min_p\lambda_{\max}(\sum_a p_aL_a)$, where
$L_a=G_a+\sum_i v_iE_{a,i}$ and $v$ is the primitive Bellman optimum.
When one attaining dual preparation $\rho_*$ satisfies
$\tr(\rho_*E_{a,i})>0$ for every command $a$ and report $i$,
the optimal-action sets $\mathcal A_i^*$ give a simpler exact memory
threshold.  Write $\mathcal P_*$ for the optimal initial rows,
$I(p)=\{i:\supp p\subseteq\mathcal A_i^*\}$, and $t(J)$ for the
least number of actions meeting all $\mathcal A_i^*$ with $i\in J$.
The inherited Theorem~\ref{thm:r77-reset-cover} states
\begin{equation}\label{eq:r77-intro-reset}
 m_*=1+\min_{p\in\mathcal P_*}
                 t(\{1,\ldots,s\}\setminus I(p)).
\end{equation}
Positive definite effects are a direct sufficient exposure hypothesis.
The explicit semidefinite example has a one-label optimum that this
cover formula would overcount without exposure.  It remains a valid
counterexample to removing that hypothesis; it is also correctly
handled by the new general frontier.  The arbitrary-cover realization
and its NP-hardness consequence are inherited unchanged.  Their
unrestricted value may be simple even though memory synthesis is hard.

\subsection*{Shared preparations and causal response realization}

Fix a finite group $G$ and finite response laws $b_i:G\to\Delta(C)$.
All preparations must use one command program and one decoder; the
decoder does not receive the preparation identity for free.  For
closed convex constraints on every command word,
Theorem~\ref{thm:r77-shared-normal-form} gives a common permutation
core, including exact charges for prescribed pure initializers.
For the distinct target tables $\mathcal B$, let $P$ range over
$G$-invariant polytopes in the full response simplex containing
$\mathcal B$, and let $v(P)$ count their vertices.  Then
Theorem~\ref{thm:r77-shared-polytope} gives
\begin{equation}\label{eq:r77-intro-polytope}
 \kappa_G(\mathcal B)=\min_P v(P),\qquad
 m_G(\mathcal B)=\min_P\{v(P)+|\mathcal B\setminus\operatorname{vert}P|\}.
\end{equation}
The polytope may extend beyond the target span.  These are different
resource conventions: a freely prepared mixed row and a prescribed
pure device label have different initialization costs.

Horizontal concatenation of the preparation Fourier blocks gives the
joint, potentially noncyclic response module.  With a common group
mean, its permutation multiplicity cover is the exact sufficiently
small positive-contrast memory, with an explicit simultaneous positive
decoder.  Weighted singular values give robust lower certificates.
Exact families on the fixed group $C_2$ require $2q$ shared labels
although every preparation separately needs two, or have mixed memory
two and unbounded pure initialization overhead.  Unequal means have
a nonnegative-rank obstruction.  The full finite-group normal form,
nonabelian multiplicity, quotient duality, and the displayed complete
symmetric-group mean--contrast law retain their proofs and scopes.

\subsection*{Operational approximation and the qubit boundary}

Finite causally prepared heralds followed by delayed terminal tests
give a robust nonnegative-factorization obstruction without requiring
a physical density matrix at each hidden simulator label.  The
decoder may have arbitrary stochastic test columns, but cannot see an
uncounted copy of the earlier history.  Theorem~\ref{thm:r77-operational-antichain}
first controls joint herald and test events, and only then conditions.
For the original qubit phase $z=(3+4i)/5$,
Theorem~\ref{thm:r77-unrestricted-qubit} implies
\begin{equation}\label{eq:r77-intro-arithmetic}
 \delta\le\frac{25\zeta}{128\cdot200^J J^2}
 \quad\Longrightarrow\quad N_{\rm cl}(\delta;B)\ge
 \min\{m:\tbinom m{\lfloor m/2\rfloor}\ge J\},
 \qquad\zeta=\min\{1,7B/8\}.
\end{equation}
At each fixed $B>0$ the unrestricted lower is of order
$\log\log(1/\delta)$; the retained constructive upper is logarithmic.
The matching leading logarithm for physical density dictionaries keeps
its Born-decoder and badly-approximable-phase hypotheses.  These rates
concern stopped response approximation, not the original general-$M$
squared-sign decision-risk frontier.  That separate problem remains
open and is not replaced by the exact reset cost hierarchy.

\subsection*{Classical ingredients, priority and the proof architecture}

Bellman verification and stochastic-shortest-path theory are classical
\cite{BertsekasTsitsiklis91}.  Finite stochastic controller search,
bounded policy iteration and nonlinear optimization of fixed controller
size have an established literature
\cite{Meuleau99,PoupartBoutilier03,AmatoBernsteinZilberstein10}.
Neither writing a nonlinear program nor invoking real algebraic
elimination is claimed as a new discovery.  The specific structural
content here is simultaneous domination with the same counted rows,
removal of all stochastic routing for the all-state reset problem,
the $s$-potential compact certificate in the undiscounted model,
and the full exact resource law and stable hierarchy.

Residual automata, positive realization, representation theory,
equivariant factorizations and antichains retain their classical
attribution \cite{RabinScott59,Rabin63,Heller65,CarlylePaz71,
BenvenutiFarina04,Schwarz64,Flor69,Serre77,GPT13,VandaeleGillisGlineur17}.
Set-cover hardness is used through its explicit physical reduction
\cite{Karp72}.  Targeted primary-source comparisons distinguish these
ingredients from the theorem statements; they do not certify
exhaustive priority or independent editorial significance.

The reset frontier depends on the experiment semantics and the finite
reset existence theorem.  It does not require the group realization or
qubit arithmetic results.  The shared-preparation chain uses the
finite-group convex normal form and its Fourier premises.  The
approximation chain uses operational factorization and phase separation.
The general actual-history comparison and paid construction remain in
the main article; their lower witnesses use the competitor's actual
histories.  The supplement preserves full independent proofs for
finite-order quantum measurement, two-stage reset, memory and
communication, rank-constrained discrimination, performed-law
compactness, variational and dual attainment, spectral memory and local
acquired geometry.  The R58 companion remains independently compilable.
No model-specific value formula is promoted to a solution of arbitrary
causal feedback, global geometry, or the independent A1/A2 programme.

\section{Prepared kernels, acquired task laws and observers}\label{sec:data}
All hidden-state, action and report spaces are standard Borel. A fixed unknown index $\theta\in\Theta$ specifies an actual preparation and positive normalized causal kernels. Positivity means nonnegative measures, not an everywhere strictly positive density. Dependence on the entire physical past may be included in the hidden state. A legal policy sees past visible reports and its admitted randomization, but not $\theta$ or a future audit. Iteration of these kernels defines its finite-prefix laws; the consistent countable path law is their kernel extension.

Each call has a cost $c_i\ge c_*>0$. The decision to make the next call is predictable. In the geometric route, the declared productive-call flag is also predictable; in the primitive-instrument route, productive reports are specified outcomes and can be random. Commitment is permitted only after a productive report, and occurs exactly once almost surely in each model. At that time a specified future audit $Z\in H$ is scored, where $H\subset\R^k$ is nonempty, compact and convex. The prediction is committed before that audit is observed. A prospective joint report--audit kernel is defined at every productive call, including rejected calls. The stopping coin is independent of the audit conditional on the entire current visible report and the past. Write $L=\operatorname{diam}(H)^2$; the case $L=0$ is trivial whenever a legal rule exists.

The core cost $C_\tau$ includes initial preparation charged to the core, failed calls and all productive calls through commitment. A separately executed final audit must be added to the raw budget. This convention is fixed before taking a risk infimum. Physical time, stored program constants and temporary computation are not identified with call count.

\subsection*{Three resource interfaces}
A checkpoint procedure may use its whole precommit history, but must output one element of a fixed dictionary of at most $M$ values. This counts its output width only; it is not a claim of $\log_2 M$ total working bits.

An autonomous observer has an alphabet $S$, $|S|\le M$, and one fixed parameter-independent program. At an intercall boundary its only persistent observer memory is its label $s\in S$. Its next command distribution is a fixed function of $s$. A physical port name, elapsed time, subprogram index, random seed or former report is not an extra free input. During a call its current command, current report and fresh coins are available. Whatever persists to another call must be encoded in $S$. A fixed decoder $r:S\to H$ reads only the committed label.

In the \emph{event-stopped} interface a post-report kernel emits $(s',e)\in S\times\{0,1\}$. An event $e=1$ ends the experiment and retains $s'$ for the decoder; no subsequent command is executed. This public commitment event is not a stored bit used in continuing computation. The same label may be used live on one path and terminal on another. In the \emph{state-stopped} interface a fixed map $h_*:S\to\{0,1\}$ determines halt status. The initial label is live, and type-one labels are globally absorbing terminal labels with no executed command rows. Entry into a terminal label is allowed only after a productive report. Thus live and terminal types are disjoint in this interface, and only in this interface by definition.

Write $R_i^\Theta(N,M)$, $i\in\{\cp,\ev,\st\}$, for the infimum of worst-case expected squared loss over programs with robust almost-sure commitment and $\sup_\theta\E_\theta C_\tau\le N$. An empty legal class has value $+\infty$. A program chosen once at random is not a free persistent seed. The ideal measurable class may have fixed real constants; effective implementation is stated separately in Section~\ref{sec:composition}.

\subsection*{A physically acquired comparison law}
Execute a declared prepared comparison experiment and obtain its joint report--audit probability $K_0(dx,dz)$ on $X\times H$. Its report marginal is $\nu$. The law is genuinely preparable and its specification has its own resource account. It need not be a member of $\Theta$. In particular no finite set of calibration data is silently promoted to a uniform certificate. Set
\begin{equation}\label{eq:derived}
 f(x)=\int zK_0(dz\mid x),\quad
 w(x)=\int\|z-f(x)\|^2K_0(dz\mid x),\quad
 d_r(x)=w(x)+\|f(x)-r\|^2.
\end{equation}
Disintegration gives Borel versions with $f\in H$; null sets may be completed arbitrarily. The profile is
\begin{equation}\label{eq:profile}
 \Q_m(b)=\min_{\substack{C\subset H,\ 1\le|C|\le m\\
                  0\le h\le b,\ \int h\,d\nu=1}}
        \int [w+\dist(f,C)^2]h\,d\nu,
 \qquad b\ge1,\ m\ge1.
\end{equation}
We use $+\infty$ for an infeasible cap or width. The quantities $f,w$ are task coordinates of acquired predictive information, not a complete substitute for the predictive quotient. Their derivation and the quotient construction in Appendix~\ref{sec:quotient} precede any use of local dimension.

\subsection*{Lower certificates on actual histories}\label{sec:lower-cert}
Fix \emph{one} actual model $\theta_*\in\Theta$, independent of the competing program. At each actual reached productive pre-call history $h$ and legal command $a$, write $P_{h,a}(dy,dz)$ for its conditional full current report--audit law. Let $P^0_{h,a}$ be a measurable comparison row on the same report space, with a Borel mark $x=\pi_{h,a}(y)$ of law $\nu$ and
\begin{equation}\label{eq:full-moments}
 \E^0[Z\mid y]=f(\pi_{h,a}y),\qquad
 \E^0[\|Z\|^2\mid y]=w(\pi_{h,a}y)+\|f(\pi_{h,a}y)\|^2.
\end{equation}
These are conditional identities given the \emph{full} report, not just its mark. Versions and inequalities are required on actual reachable histories, not only almost everywhere under an unrelated comparison path law.

Fix $0<l_-\le1\le u_-<\infty$. For all Borel $q:Y\to[0,1]$ and all finite-image Borel $r:Y\to H$, suppose predictable nonnegative defects $a_-(h,a),e_-(h,a)$ obey
\begin{align}
 \int q\,d(P_{h,a})_Y
 &\le u_-\int q\,d(P^0_{h,a})_Y+a_-(h,a),\label{eq:lower-mass}\\
 \int q(y)\|z-r(y)\|^2P_{h,a}(dy,dz)
 &\ge l_-\int q(y)\|z-r(y)\|^2P^0_{h,a}(dy,dz)-e_-(h,a).
 \label{eq:lower-task}
\end{align}
Independent current randomization can be adjoined to both rows: integrating inequalities for each fixed random value proves the same inequalities with the same defects. All command-conditioned lower inequalities hold before averaging the chosen command.

Let $A_\tau$ count productive calls. At this same actual model require, for every legal competing procedure,
\begin{equation}\label{eq:effort}
 \E_* C_\tau\ge\kappa_-\E_* A_\tau,\qquad\kappa_->0.
\end{equation}
At budget $N$ suppose, uniformly over all robustly feasible competitors,
\begin{equation}\label{eq:occupation}
 \E_*\sum_{j\le\tau:\ j\ {\rm productive}}a_-(H_{j-1},a_j)\le A_N,
 \qquad
 \E_*\sum_{j\le\tau:\ j\ {\rm productive}}e_-(H_{j-1},a_j)\le E_N.
\end{equation}
These are actual occupation certificates, not properties merely of the proposed upper program. Uniform row defects give $A_N\le(N/\kappa_-)\bar a_-$ and $E_N\le(N/\kappa_-)\bar e_-$. Lemma~\ref{lem:drift} supplies a nonuniform alternative.

\subsection*{One common upper construction}\label{sec:upper-cert}
Independently, let $\mathcal G$ be a parameter-independent program which, from each of its start histories, reaches one productive report. Its conditional expected attempt cost is at most $\kappa_+>0$ on every surviving history in every $\theta\in\Theta$. It uses $B$ live roles and returns to its start label on rejection; no reset of the physical state is assumed. It forms a virtual mark $x$, a current choice $J$ and a permitted export $T_Jr$ of its task readout. For every finite dictionary $C$ assume
\begin{equation}\label{eq:exports}
 |\{T_jr:j\in J,\ r\in C\}|\le D_0|C|,
 \qquad D_0\in\mathbb N.
\end{equation}
Every continuing phase and coin needed for this construction is included in the $B$ live roles. For fixed $0<l_+\le1\le u_+<\infty$ and $a_+,e_+\ge0$, assume for all Borel $q:X\to[0,1]$ and finite-image Borel $r:X\to H$, conditionally on every such history,
\begin{align}
 \E[q(x)]&\ge l_+\int q\,d\nu-a_+,\label{eq:upper-mass}\\
 \E[q(x)\|Z-T_Jr(x)\|^2]
 &\le u_+\int q(x)d_{r(x)}(x)\nu(dx)+e_+.
 \label{eq:upper-task}
\end{align}
The expectation here includes the within-attempt choice $J$. Unlike the lower, the upper is not required to hold for each choice separately. The choices retained after this call are nevertheless charged through the actual export or live role. These hypotheses express one common causal construction, not modelwise optimizers.

\section{Primitive optimal control for finite reset instruments}\label{sec:reset-control}
An exact application should not require that a balancing group has already been built into the experiment. We give a class in which the best continuation and a least-favorable preparation can instead be derived from the primitive measurement instrument. The class includes any finite family of rank-one failure projections in any finite Hilbert dimension, as well as mixed post-failure states. It does not subsume an arbitrary-rank failure projection, whose posterior still depends on its input. The complete arbitrary-rank two-detector theorem is retained in Section~\ref{sec:two-detectors} of the supplement.

Let $\mathcal H$ be finite-dimensional. There is a nonempty finite command set $\mathcal A$, nonterminal reports $i=1,\ldots,s$, and a terminal success. Assume $s\ge1$. A nonterminal branch has the physically specified map
\begin{equation}\label{eq:reset-instrument}
 \mathcal I_{a,i}(\rho)=\tr(E_{a,i}\rho)\sigma_i,
 \qquad E_{a,i}\succeq0,\quad \sum_i E_{a,i}\preccurlyeq I,
\end{equation}
where $\sigma_i$ is a known density matrix and the report $i$ is observed. The aggregate success effect is $B_a=I-\sum_iE_{a,i}$. The instrument itself prepares the state $\sigma_i$ on failure; no additional freely available reset is assumed. One call has a nonnegative cost with conditional mean $\tr(G_a\rho)$, where
\[
 0<c_*I\preccurlyeq G_a\preccurlyeq c^*I<\infty.
\]
Success requires immediate commitment and failure forbids it. The decoder has one fixed terminal value. All commands are legal. The initial state is arbitrary and unknown to the observer; the known post-report state is part of the physical model, not an unknown-state oracle. Commands and reports may be relabelled without changing these data.

Set
\[
 Q_{iaj}=\tr(E_{a,j}\sigma_i),\qquad c_{ia}=\tr(G_a\sigma_i),
 \qquad q_{ia}^{\rm suc}=1-\sum_jQ_{iaj}.
\]
A nonempty set $S\subset\{1,\ldots,s\}$ is a \emph{closed failure set} when $\sum_{j\in S}Q_{iaj}=1$ for every $i\in S$ and every command $a$.

\begin{theorem}[Finite reset classification, optimum and preparation saddle]\label{thm:reset-control}
There exists one common visible-history policy which is proper at finite uniform expected cost for every initial density matrix if and only if there is no closed failure set. If there is none, the finite Bellman system
\begin{equation}\label{eq:reset-bellman}
 v_i=\min_{a\in\mathcal A}\left(c_{ia}+\sum_jQ_{iaj}v_j\right),\qquad 1\le i\le s,
\end{equation}
has a unique finite nonnegative solution. It is the optimal remaining cost from each reported state $\sigma_i$. A deterministic minimizing command at every reported state attains it and is proper.

Define Hermitian start-cost operators
\[
 L_a=G_a+\sum_jv_jE_{a,j}.
\]
The unknown-initial-state value is exactly
\begin{align}
 V_*&=\min_{p\in\Delta(\mathcal A)}\lambda_{\max}\left(\sum_a p_a L_a\right)\label{eq:reset-sdp}\\
 &=\max_{\rho\succeq0,\ \tr\rho=1}\min_a\tr(\rho L_a).\label{eq:reset-dual}
\end{align}
Both extrema are attained. A minimizing initial command distribution followed by the deterministic Bellman continuation is one robustly minimax policy whose uniform bound holds for every initial state; its initial row need not minimize cost separately at each known state. A maximizing density matrix is a primitive least-favorable preparation against every visible-history policy. The latter statement is for unrestricted history; it is not automatically a saddle at every smaller memory budget.

If $g$ distinct command values occur in a chosen deterministic Bellman selector, then $g+1\le\min(s,|\mathcal A|)+1$ live labels suffice with a prescribed initial label. If a minimizing initial distribution is a pure command already used by that selector, $g$ suffice. Event stopping permits reuse of a live label for the constant terminal value; state stopping adds one disjoint terminal label.
\end{theorem}
\begin{proof}
If a closed failure set exists, start in any $\sigma_i$ with $i\in S$. Every command fails with probability one and reports another state in $S$. Induction makes success impossible under every policy; the positive cost floor makes its expected cost infinite.

Suppose there is no such set. Choose every command independently with strictly positive probability. From each failure state there is a positive-probability directed path to success: otherwise the set of failure states reachable from that state would be nonempty and closed under every command, contradicting the assumption. A shortest such path visits no failure state twice and has at most $s$ edges. Finiteness gives a positive lower bound $\delta$ for the probability of following one such path in the next $s$ calls, uniformly in its starting failure state. Consequently survival after $ns$ calls is at most $(1-\delta)^n$. From an arbitrary initial state the first call either succeeds or enters one of the reported states. The expected number of calls is therefore uniformly at most $1+s/\delta$, and the expected cost at most $c^*(1+s/\delta)$. This proves existence using even a one-live-label randomized command row. Its role is feasibility, not optimality.

For explicit primitive constants, if the row chooses each command with probability at least $\eta>0$ and $\epsilon>0$ is the smallest positive number among the finite collection $Q_{iaj},q_{ia}^{\rm suc}$, then one may take $\delta=(\eta\epsilon)^s$. Factors for a shorter path can be bounded below by this same number. Zero probabilities are not divided out or assigned a false positive floor.

Start value iteration at $v^{(0)}=0$. The finite-horizon dynamic programming recursion gives $v^{(n+1)}_i=\min_a(c_{ia}+\sum_jQ_{iaj}v^{(n)}_j)$. These values increase and are bounded by the proper randomized policy's finite cost. Their limit $v$ is finite and solves \eqref{eq:reset-bellman}, since a finite minimum commutes with this coordinatewise limit. Select a minimizing command $\alpha(i)$ in each row. Iterating its equality bounds every partial expected cost by $v_i$. The positive floor then bounds the sum of survival probabilities, so the selected policy is proper and has finite cost. The bounded terminal remainder $\E[v_{I_n}\one_{\{\tau>n\}}]$ vanishes. Passing to the limit in the iterated equality identifies its cost as $v_i$.

For any other proper policy, Bellman's inequality for each chosen command telescopes to $v_i$ no greater than its expected cost plus the same bounded survival remainder. Letting the horizon grow proves its cost at least $v_i$. An infinite-cost policy cannot improve this bound. The same argument for any finite nonnegative Bellman solution $w$ shows that its greedy policy realizes $w$ and that $w$ is below the cost of every proper policy. Comparing to $v$ in both directions proves uniqueness.

Before the first call the observer has no report, so its first command has one model-independent distribution $p$. After a failure $j$, the physical state is exactly $\sigma_j$, irrespective of the original state or previous private coins. The already proved continuation lower therefore makes its initial cost at least
$\tr[\rho\sum_a p_aL_a]$. The Bellman continuation attains this bound for every $\rho$. Taking the worst density matrix gives the largest eigenvalue, proving \eqref{eq:reset-sdp}. Compactness of the probability simplex gives its minimum.

For completeness the dual equality follows directly from finite-dimensional separation. Let $\mathcal K=\operatorname{conv}\{L_a\}$ and let $v_* =\min_{L\in\mathcal K}\lambda_{\max}(L)$. The compact convex set $\mathcal K$ is disjoint from the open convex set $\{H:H\prec v_*I\}$. A separating real linear functional on Hermitian matrices is pairing with a nonzero Hermitian $R$, with $\inf_{L\in\mathcal K}\tr(RL)\ge\sup_{H\prec v_*I}\tr(RH)$. Finiteness of the right side forces $R\succeq0$, since otherwise moving $H$ arbitrarily far in a negative semidefinite direction makes that pairing unbounded above. Normalize $\tr R=1$; the right supremum is then $v_*$. This gives a density matrix with $\min_a\tr(RL_a)\ge v_*$. Conversely, pairing any density matrix with a minimizing $L$ bounds that minimum by $v_*$. This proves equality and dual attainment. The continuation lower makes the maximizing state a witness against every policy, not merely against first-call mixtures.

To implement the policy, use one label for its initial distribution and a pure command label for each distinct value of $\alpha$. On failure report $j$, retain the label for $\alpha(j)$. The next row depends only on that retained label. Success commits immediately. If the initial row is already one of these pure rows, use that existing label as the start. The event/state packing lemma gives the terminal counts. No running clock, earlier report or persistent seed is available in addition to the labels.
\end{proof}

\begin{proposition}[Equality, rigidity and a memory lower certificate]\label{prop:reset-rigidity}
Let $\rho_*$ be an attaining preparation in \eqref{eq:reset-dual}, and put
\[
 d_{ia}=c_{ia}+\sum_jQ_{iaj}v_j-v_i\ge0,\qquad
 \mathcal A_i^*=\{a:d_{ia}=0\}.
\]
For every finite-cost policy with first-command distribution $p$, its cost satisfies the exact residual identity
\begin{equation}\label{eq:reset-residual}
 \E_{\rho_*}C
 =\tr\!\left(\rho_*\sum_a p_aL_a\right)
    +\E_{\rho_*}\sum_{n\ge1:\,\tau>n}d_{I_n,A_{n+1}}.
\end{equation}
Consequently an unrestricted optimum uses only commands in $\mathcal A_i^*$ at every positive-probability continuation from reported state $i$, under this preparation.

Suppose a set $I$ of reported states is visited with positive probability by every robust optimum evaluated at $\rho_*$. Let $\chi(I)$ be the least number of subsets covering $I$ whose respective intersections of optimal command sets are nonempty. Every fixed-row autonomous optimum needs at least $\chi(I)$ live labels. If every minimizing initial row has support not contained in $\mathcal A_i^*$ for every $i\in I$, it needs at least $\chi(I)+1$. These are checkable lower certificates once the mandatory set and the initial-row condition have been proved, not assertions that every experiment has the same minimal memory.
\end{proposition}
\begin{proof}
After each failed call, subtract $v_{I_n}$ from the cost of the next call plus the conditional next potential. The conditional difference is exactly $d_{I_n,A_{n+1}}$. Sum to a deterministic horizon, including the first-call cost and potential. The remaining potential has absolute value at most $\max_i v_i$ times the survival indicator. Finite expected cost and the floor make survival tend to zero. Monotone convergence for the nonnegative residual sum yields \eqref{eq:reset-residual}. Its first term is at least $V_*$ by the dual witness. A robust optimal policy's cost at $\rho_*$ is at most $V_*$, so every visited positive-probability residual vanishes.

A retained label serving states in a subset $J$ has one command row. Its support must lie in every $\mathcal A_i^*$, $i\in J$, because each such visit has positive probability and any action of positive row weight with positive residual would contradict equality. Thus $\bigcap_{i\in J}\mathcal A_i^*$ is nonempty. The labels used by mandatory states cover $I$, proving the first lower. For an optimum, its first row is a minimizing row in \eqref{eq:reset-sdp}: its worst-state cost dominates the first-call lower, whose minimum is $V_*$. Under the additional support condition this initial row can serve none of the mandatory states, so its label is additional to their covering labels. This proves the sharpened count.
\end{proof}

\begin{corollary}[Arbitrary finite rank-one detector families]\label{cor:rank-one-family}
Let the failure projection of detector $a$ be $P_a=|u_a\rangle\langle u_a|$ on any finite-dimensional Hilbert space, with $\|u_a\|=1$, and let its deterministic call cost be $k_a>0$. There is no intervening evolution or additional control. If all rays $\mathbb C u_a$ coincide, that ray is dark and the all-state value is infinite. Otherwise the value is finite and is given by Theorem~\ref{thm:reset-control}, with
\[
 v_i=\min_a\{k_a+|\langle u_i,u_a\rangle|^2v_a\},
 \qquad L_a=k_aI+v_aP_a.
\]
Thus the continuation is a finite scalar system, while the unknown-initial-state optimization retains the full spectral geometry of the physical projections. There is no fixed prior or restriction to two detectors or two physical dimensions.
\end{corollary}
\begin{proof}
A rank-one failure satisfies $P_a\rho P_a=\tr(P_a\rho)P_a$, so it is a reset branch with report $a$ and state $P_a$. If all rays agree, their common ray never succeeds. Otherwise choose two distinct rays and alternate their detectors. Their squared overlap is less than one. After the first failure every switch has this same strictly positive success probability, giving finite uniform cost. The classification theorem applies. Substitution of the reset transition probabilities gives the displayed formulas. The general equality and memory certificates retain the actual observed failure label and therefore apply to this family as well.
\end{proof}

The reduction to a finite absorbing decision problem, Bellman verification, convex separation and spectral duality are classical methods; quantum decision-process frameworks already provide much broader control languages \cite{Barry14,Ying18,Saldi25}. The contribution claimed here is the explicit primitive reset-instrument criterion, the common all-initial-state spectral value, its realization with charged rows, and the residual-based memory certificate within this experiment. It is distinct from optimizing the probing interval for a fixed detector with unitary evolution \cite{Kessler21,DelVecchio26}. It also does not replace the preserved two-stage reset communication and rank-constrained discrimination results by a first-detection theorem.

\section{The complete finite-memory frontier for reset control}
\label{sec:r78-reset-frontier}

The optimal-memory threshold and the cost at every smaller memory are
different invariants.  We now determine the latter by an attained finite
variational principle.  The result applies to the full reset class of
Section~\ref{sec:reset-control}; strict positivity of the failure effects
is not required.  In particular it includes reports which a least-favorable
preparation cannot reach.  The key reduction preserves the controller's
existing command rows while removing all dependence of the post-report
update on the preceding label and command.

Keep the primitive data $E_{a,j},G_a,\sigma_i$, the positive cost floor,
and the absence of a closed failure set from
Theorem~\ref{thm:reset-control}.  For a command distribution
$p\in\Delta(\mathcal A)$ define
\[
 c_i(p)=\sum_a p_a c_{ia},\qquad
 Q_{ij}(p)=\sum_a p_a Q_{iaj},\qquad
 G(p)=\sum_a p_aG_a,\qquad E_j(p)=\sum_a p_aE_{a,j}.
\]
All these quantities concern the actual reported reset state.  There is
no additional measurement, cost observation, free clock, or access to the
unknown initial density matrix.

Let $V_M$ be the infimum of the worst-initial-state expected cost over
autonomous controllers with at most $M\ge1$ live labels, a prescribed
pure initial label, and mandatory stopping at success.  The command row
at a label can be mixed.  Updates after a failure may initially depend
stochastically on the old label, the sampled command, and its report.
Controllers with infinite worst-state cost do not affect the infimum.

\subsection{A memory-preserving deterministic update theorem}

\begin{theorem}[Dominance with the same command rows]
\label{thm:r78-palette-normal-form}
Suppose an $M$-label controller has finite worst-initial-state cost $C$,
and its command rows are $p^1,\ldots,p^M$, with pure initial label $1$.
There is another controller with these same rows and the same initial
label whose update after report $i$ is one deterministic label $k(i)$,
independent of the preceding label and command.  Its expected cost is
no greater than the original cost at every initial density matrix.

More precisely, the finite system
\begin{equation}\label{eq:r78-palette-bellman}
 x_i=\min_{1\le k\le M}
       \left\{c_i(p^k)+\sum_jQ_{ij}(p^k)x_j\right\}
\end{equation}
has a unique finite nonnegative solution, with $0\le x_i\le C$.
Choose $k(i)$ attaining the minimum.  The resulting continuation is
proper from every reset state, and its full cost from an initial $\rho$
is exactly
\begin{equation}\label{eq:r78-palette-start}
 \tr\bigl(\rho H(p^1,x)\bigr),\qquad
 H(p,x)=G(p)+\sum_jx_jE_j(p).
\end{equation}
This is a pointwise cost comparison, without a convexification of the
hard memory constraint.
\end{theorem}
\begin{proof}
Regard the finite list of rows as the available choices at each reset
state.  A choice $k$ has mean cost $c_i(p^k)$ and substochastic reset
transition row $Q_{ij}(p^k)$.  Any extra information about the sampled
command or earlier private coins does not alter the physical state
$\sigma_j$ after report $j$.

The original controller can be started with label $1$ and physical
state $\sigma_i$, for every $i$: the robustness quantifier includes
\emph{all} initial density matrices.  This gives, for the same list
of rows, a proper history-dependent policy of cost at most $C$ from
each reset state.  This comparison restart is used only to prove
finiteness; it is not an additional operation supplied to the eventual
controller.  In particular it remains valid when some reset reports
are inaccessible from a least-favorable preparation.

Starting the finite-horizon recursion for the displayed row choices
at zero gives increasing vectors bounded above by $C\mathbf1$.
Their limit $x$ is finite and satisfies
\eqref{eq:r78-palette-bellman}.  For a minimizing selector $k(i)$,
iterate its equality through a finite horizon.  The accumulated
expected cost is bounded by $x_i$.  The positive cost floor therefore
bounds the sum of survival probabilities, so the selector is proper
from every state.  The bounded remaining potential tends to zero;
its expected total cost is $x_i$.

Conversely Bellman's inequality holds conditionally at each call of
any policy using these rows, even if it retains the sampled command
and earlier coins.  Telescope it to a finite horizon.  For a
finite-cost policy, the survival remainder vanishes by the same
positive floor, proving that its cost from reset state $i$ is at
least $x_i$.  Infinite-cost policies cannot improve the bound.  The
same verification applied to any finite nonnegative solution gives
uniqueness.

After the original first call, each reached continuation from report
$j$ uses the same list of rows and hence costs at least $x_j$.
Conditioning on that call proves that its cost at every $\rho$ is
at least \eqref{eq:r78-palette-start}.  The new controller starts
with row $p^1$ and, after report $j$, writes precisely the existing
label $k(j)$.  Thereafter it implements the minimizing selector and
attains this expression.  The actual device thus uses at most the
same $M$ labels, with no uncounted row-selection memory.
\end{proof}

The all-initial-state quantifier is essential to the bounded comparison
in this proof.  Properness for just one supplied initial state would
not by itself imply a finite continuation value at every reset state.
Nor does the theorem assert that mixed command rows can always be
replaced by pure commands: only their post-report selection is made
deterministic.

\subsection{A compact exact variational principle}

Choose once and for all a full-support row $\bar p$, for example the
uniform row on $\mathcal A$.  The proof of
Theorem~\ref{thm:reset-control} makes $\bar Q=(Q_{ij}(\bar p))$
transient.  Consequently
\begin{equation}\label{eq:r78-explicit-cap}
 \bar x=(I-\bar Q)^{-1}(c_i(\bar p))_{i=1}^s,
 \qquad C_0=\lambda_{\max}H(\bar p,\bar x)<\infty.
\end{equation}
The constant one-row controller has worst-state cost $C_0$, and
$\bar x_i=\tr(\sigma_iH(\bar p,\bar x))\le C_0$.
Thus this is a primitive, explicit compactness bound, not an assumed
upper bound on an unknown optimal trajectory.

For $k:\{1,\ldots,s\}\to\{1,\ldots,M\}$ let $\mathcal F_{M,k}$
consist of the rows $p^1,\ldots,p^M\in\Delta(\mathcal A)$ and
$x\in[0,C_0]^s$ satisfying
\begin{equation}\label{eq:r78-supersolution}
 x_i\ge c_i(p^{k(i)})+\sum_jQ_{ij}(p^{k(i)})x_j
 \qquad(1\le i\le s).
\end{equation}
There are $M^s$ possible selectors and each corresponding feasible
set is compact; some of them may be empty.

\begin{theorem}[Exact attained frontier at every memory]
\label{thm:r78-reset-frontier}
For every $M\ge1$,
\begin{equation}\label{eq:r78-full-frontier}
 V_M=\min_k\ \min_{(p^1,\ldots,p^M,x)\in\mathcal F_{M,k}}
             \lambda_{\max}H(p^1,x).
\end{equation}
The minimum is finite and attained, and an attaining controller has
deterministic report-dependent updates as in
Theorem~\ref{thm:r78-palette-normal-form}.  For $0\le b\le C_0$,
there is an $M$-label controller of worst-state cost at most $b$ if
and only if one of the finite systems
\begin{equation}\label{eq:r78-budget-certificate}
 (p^1,\ldots,p^M,x)\in\mathcal F_{M,k},\qquad
 H(p^1,x)\preccurlyeq bI
\end{equation}
is feasible.

In particular $V_M$ is nonincreasing in $M$ and equals the unrestricted
value $V_*$ once $M\ge\min(s,|\mathcal A|)+1$.  The exact optimal
label count, without a strict-exposure hypothesis, is
\begin{equation}\label{eq:r78-general-threshold}
 m_{\mathrm{opt}}=
 \min\{M:\text{\eqref{eq:r78-budget-certificate} is feasible at }b=V_*\}.
\end{equation}
Every $M<m_{\mathrm{opt}}$ has the strict gap $V_M>V_*$.
For algebraic primitive data, $V_M$, $m_{\mathrm{opt}}$, and an
attaining controller admit finite exact real-algebraic procedures.
No polynomial complexity or general minimax saddle at fixed $M$ is
asserted.
\end{theorem}
\begin{proof}
A feasible supersolution bounds every finite-prefix continuation
cost under the displayed selector by $x_i$.  The positive cost floor
therefore makes that selector proper from every reset state.  Its
true continuation vector is at most $x$, and its full expected cost
operator is at most $H(p^1,x)$, because every $E_j(p^1)$ is positive
semidefinite.  Thus every feasible point is an actual controller
upper certificate.

Conversely take any controller of worst-state cost $b\le C_0$.
Theorem~\ref{thm:r78-palette-normal-form} supplies its row-list
Bellman vector $x\le b\mathbf1$, a minimizing selector, and a
cost operator bounded by $bI$.  These are feasible in
\eqref{eq:r78-budget-certificate}.  The bound $x\le b\mathbf1$
uses the comparison starting at each $\sigma_i$ with label $1$;
it is not inferred from positive exposure of a particular dual state.

All rows equal to $\bar p$, with $x=\bar x$, give a feasible point
of value $C_0$, for every $M$.  A finite union of the compact
$\mathcal F_{M,k}$ is therefore nonempty and the continuous
objective attains its minimum.  If the controller infimum is strictly below $C_0$, apply the converse
to an approximating sequence below $C_0$; if it equals $C_0$, use
the displayed one-row certificate.  Together with the first implication,
this identifies the minimum with the controller infimum.  At a minimizing feasible
point the constructed controller cannot have strictly smaller
worst-state cost, since that would contradict the equality of the
infima.  This proves actual attainment as well as the budget
equivalence.

Allowing duplicate rows shows monotonicity.  The unrestricted
implementation in Theorem~\ref{thm:reset-control} proves the stated
finite saturation bound.  At every smaller $M$, equality with $V_*$
is equivalent to an attaining optimal controller by the result
already proved, giving \eqref{eq:r78-general-threshold} and the
strict-gap assertion.  This supplies a general variational
classification even where the simpler transversal formula of
Theorem~\ref{thm:r77-reset-cover} is inapplicable.

For algebraic data, $\bar x$ and $C_0$ are algebraic.  For each
selector, the row constraints and
\eqref{eq:r78-supersolution} are polynomial inequalities with
algebraic coefficients.  On separating real and imaginary parts,
the Hermitian inequality in \eqref{eq:r78-budget-certificate} is
equivalent to nonnegativity of its principal minors.  Exact
real-algebraic elimination determines its feasible budget set and
least endpoint; compactness ensures that endpoint is included and
provides algebraic attaining coordinates.  Finally enumerate the
finite range in \eqref{eq:r78-general-threshold}.  This is an exact
computability conclusion from the proved finite system, not a
claim that a numerical optimization certifies a tie or that an
elimination computation has been performed for arbitrary inputs.
\end{proof}

With the inherited single terminal decoder, event stopping has the
same total-label frontier $V_M$.  A state-stopped device with $M+1$
total labels has the live-label frontier $V_M$, since its terminal
label is disjoint.  The prescribed pure start is already charged
in \eqref{eq:r78-full-frontier}: $p^1$ is the command distribution
at a counted initial label, not a freely mixed choice of hidden
initial labels.

\begin{proposition}[Continuity of the full frontier]
\label{prop:r78-frontier-stability}
Fix the Hilbert dimension and the finite command and report alphabets.
For every fixed $M$, $V_M$ is continuous in the primitive effects,
costs, and reset states on the set of valid reset instruments with a
positive cost floor and no closed failure set.  Thus if $V_M>V_*$ at
one instrument, the same strict inequality holds throughout some
neighborhood within this class.  This assertion concerns values and
strict gaps; the exact optimal-memory count need not be continuous.
\end{proposition}
\begin{proof}
Near a fixed instrument, the same full-support row $\bar p$ has
transient $\bar Q$: its spectral radius is strictly less than one,
an open condition.  The resolvent in
\eqref{eq:r78-explicit-cap} and its spectral value $C_0$ are
continuous, so the compactness bounds are locally uniform.

An attaining controller can be chosen with deterministic report
updates and with its true continuation vector finite from every
reset state.  Its substochastic reset matrix $Q^k$, whose $i$th
row is $Q_{ij}(p^{k(i)})$, is transient.  Indeed its finite expected
cost and the positive cost floor bound expected duration from every
state, so the entries of $\sum_{n\ge0}(Q^k)^n$ are finite.
Keep its rows and selector fixed under a small perturbation.  The
same matrix remains transient, and both
$(I-Q^k)^{-1}(c_i(p^{k(i)}))_i$ and the resulting initial cost
operator vary continuously.  This proves upper semicontinuity of
$V_M$.

For the reverse inequality, take a convergent sequence of primitive
data and attaining points in \eqref{eq:r78-full-frontier}, along
a subsequence realizing the lower limit of the values.  The locally
uniform cap bounds the continuation vectors, the rows lie in compact
simplexes, and a further subsequence has a constant selector.  Pass
to a limit in the supersolution inequalities and in the initial
matrix upper bound at their attaining values.  The limit is an
actual controller certificate by the first part of the preceding
proof.  Its value is no greater than the lower limit, proving lower
semicontinuity.  This argument does not need the limiting continuation
vector to lie below one particular choice of cap; its finite
supersolution already constructs a proper controller.

The unrestricted value equals $V_{\min(s,|\mathcal A|)+1}$, so it
is continuous as well.  Continuity of their difference proves the
strict-gap statement.
\end{proof}

The domination step is specific to observed resetting: after each
report the physical state no longer depends on the earlier record.
Bellman verification, finite-state-controller optimization, spectral
minimax at unrestricted memory, and real-algebraic elimination are
classical ingredients
\cite{BertsekasTsitsiklis91,Meuleau99,PoupartBoutilier03,AmatoBernsteinZilberstein10,Seidenberg54}.
The present result combines them in an exact
counted-memory normal form and compact necessary-and-sufficient
budget certificate.  It does not prove such a reduction for a
general partially observed process or for the original squared-sign
risk frontier.

\subsection{Closed frontiers for arbitrary action-cover instances}

The preceding variational formula can be solved explicitly for the
rational family in Theorem~\ref{thm:r77-reset-universality}.  Let
$A_1,\ldots,A_s$ be nonempty proper subsets of $[q]=\{1,\ldots,q\}$,
$q\ge2$.  On the basis $e_1,\ldots,e_s,f_1,\ldots,f_q$, keep
\[
 \sigma_i=|e_i\rangle\langle e_i|,\quad E_{a,i}=\frac rs I,
 \quad G_a(e_i)=1+\mathbf1_{\{a\notin A_i\}},\quad
 G_a(f_j)=3+\mathbf1_{\{a=j\}},\qquad 0\le r<1.
\]
Here $G_a(e_i)$ and $G_a(f_j)$ denote the corresponding diagonal
entries.  The success effect is $(1-r)I$.  This also defines the
degenerate endpoint $r=0$, where every first call succeeds.  For
$H\subseteq[q]$ put
\[
 U(H)=\{i:H\cap A_i=\varnothing\},\qquad
 b_a(H)=|\{i\in U(H):a\notin A_i\}|,
 \qquad \lambda=\frac{r}{s(1-r)}.
\]
Let $b_{(1)}(H)\le\cdots\le b_{(q)}(H)$ be these $q$ nonnegative
integers in increasing order, with repetitions.

\begin{theorem}[The complete action-cover cost frontier]
\label{thm:r78-cover-frontier}
For every $M\ge1$, the exact worst-state cost at $M$ live labels is
\begin{equation}\label{eq:r78-cover-spectrum}
 V_M=3+\frac r{1-r}+
 \min_{\substack{H\subseteq[q]\\ |H|\le M-1}}
 \ \min_{1\le k\le q}
       \frac{1+\lambda\sum_{j=1}^k b_{(j)}(H)}{k}.
\end{equation}
An optimum uses pure continuing command rows indexed by $H$ and an
initial row uniform on $k$ commands with the $k$ smallest values of
$b_a(H)$.  The initial row and these continuing rows suffice for
deterministic report-dependent routing.  Both minima are finite
enumerations, including all tied and degenerate cases.

For a fixed set system, after subtracting $3+r/(1-r)$ the frontier
is a continuous, nondecreasing, piecewise affine concave function
of $\lambda\ge0$.  At $r=0$ all memory values equal $3+1/q$ and
one label is optimal; for $r>0$ the unrestricted-memory threshold
is the transversal number plus one, as in
Theorem~\ref{thm:r77-reset-universality}.
\end{theorem}
\begin{proof}
Fix any list of rows and put
\[
 h_i=\min_{1\le k\le M}p^k(A_i^c),\qquad
 \overline h=\frac1s\sum_i h_i.
\]
The reset probabilities are independent of the command, so
\eqref{eq:r78-palette-bellman} has the explicit solution
\[
 x_i=1+h_i+\frac{r(1+\overline h)}{1-r},
 \qquad \frac1s\sum_i x_i=\frac{1+\overline h}{1-r}.
\]
Its first-cost eigenvalues on $f_j$ are
$3+p^1_j+r(1+\overline h)/(1-r)$; those on $e_i$ are at most
$2+r(1+\overline h)/(1-r)$.  Hence the former strictly dominate,
and the exact cost for these rows after optimal routing is
\begin{equation}\label{eq:r78-cover-row-objective}
 3+\frac r{1-r}+\max_a p^1_a+
               \lambda\sum_i\min_k p^k(A_i^c).
\end{equation}
This formula applies to all row choices; success has probability
$1-r>0$ on every call, so no properness exception is present.

For each row $p^k$ with $k\ge2$, with the other rows fixed, the
last sum in \eqref{eq:r78-cover-row-objective} is a continuous
concave function of $p^k$: each summand is the minimum of one
linear function and constants.  Its minimum over a simplex has a
vertex minimizer.  Start from a global minimum, which exists by
compactness, and replace these rows successively by such vertices.
No replacement can increase the objective; global minimality then
keeps it unchanged.  Thus all continuing rows may be pure.

Let $H$ be their distinct actions.  When $i\notin U(H)$ some
continuing row has zero failure penalty and $h_i=0$.  When
$i\in U(H)$ all continuing rows have penalty one, giving
$h_i=p^1(A_i^c)$, also when $H$ is empty.  It follows that
\[
 \sum_i h_i=\sum_a b_a(H)p^1_a.
\]
The remaining initial-row minimization is therefore
\begin{equation}\label{eq:r78-initial-cap-problem}
 \min_{p\in\Delta([q])}\left\{\max_a p_a+
                                  \lambda\sum_a b_a(H)p_a\right\}.
\end{equation}

For completeness introduce $t\in[1/q,1]$ and impose $p_a\le t$.
For fixed $t$, the linear term is minimized by filling the smallest
$b_a(H)$ coordinates first, each up to $t$, with at most one
partially filled coordinate.  An exchange of mass from a larger
coefficient to a smaller unsaturated one proves this assertion,
including ties.  On each interval
$[1/(k+1),1/k]$ the minimized expression
$t+\lambda\sum_a b_a(H)p_a$ is affine in $t$.  Its minimum has an
endpoint minimizer; the endpoints are exactly $t=1/k$,
$1\le k\le q$.  At such an endpoint a minimizing $p$ is uniform
on $k$ smallest coefficients, and its value is
$(1+\lambda\sum_{j=1}^k b_{(j)}(H))/k$.  The epigraph problem in
$t,p$ has the same optimum as \eqref{eq:r78-initial-cap-problem},
since one may always reduce $t$ to $\max p_a$.  This proves
\eqref{eq:r78-cover-spectrum} and the attaining construction.

For fixed $H,k$ the displayed value after subtracting the common
term is affine with nonnegative slope in $\lambda$.  The finite
minimum is therefore continuous, piecewise affine, concave and
nondecreasing.  At $r=0$ its minimum is $1/q$, attained by a
single uniform row.  For $r>0$, an optimal initial row at the
unrestricted value must be uniform on all $q$ actions; because
each $A_i$ is proper it has positive penalty at every uncovered
state.  Thus zero excess cost requires $H$ to meet every $A_i$,
and a smallest transversal is sufficient.  Equivalently this is
the already proved positive-effect threshold.
\end{proof}

\begin{corollary}[Arbitrarily long strict memory hierarchies]
\label{cor:r78-strict-hierarchy}
For every $n\ge2$ there is a rational reset experiment in dimension
$2n$, with $n$ commands, $n$ failure reports, success probability
$1/2$ at every call and all costs in $[1,4]$, for which
\begin{equation}\label{eq:r78-strict-hierarchy}
 V_M=4+\frac1n+
       \frac{(n-1)\max\{n-M+1,0\}}{n^2}
       \qquad(M\ge1).
\end{equation}
In particular $V_1>V_2>\cdots>V_{n+1}=V_*$, and each of the
first $n$ memory increments saves exactly $(n-1)/n^2$ in worst-state
expected cost.  All failure effects are positive definite.  In this
family the retained first-residual certificate satisfies
\begin{equation}\label{eq:r78-certificate-gap}
 D_M=\frac{V_M-V_*}{2}.
\end{equation}
Thus it is a strict lower certificate below the optimum threshold,
whereas \eqref{eq:r78-strict-hierarchy} is the full repeated-experiment
cost frontier.

The $n$ strict inequalities persist in an open neighborhood within
the class of valid instruments.  They also persist under explicit
arbitrarily small perturbations for which a failure effect and a
cost operator do not commute.  These perturbations retain positive
failure effects, dimension $2n$, and mean cost operators
$I\preccurlyeq G_a\preccurlyeq4I$; the exact
formula \eqref{eq:r78-strict-hierarchy} refers to the unperturbed
instrument.
\end{corollary}
\begin{proof}
Use the preceding construction with $q=s=n$, $A_i=\{i\}$ and
$r=1/2$, so $\lambda=1/n$.  For a continuing action set $H$,
write $u=n-|H|$ for the number of uncovered singleton sets.
Then $b_a(H)=u-1$ on the $u$ uncovered actions and $b_a(H)=u$
on the remaining actions.  If $u>0$, the candidate initial-row
value at a support size $k$ is
\[
 \begin{cases}
  k^{-1}+(u-1)/n,&1\le k\le u,\\[2pt]
  u/n+(1-u/n)/k,&u\le k\le n.
 \end{cases}
\]
The first expression is minimized at $k=u$; the second is
nonincreasing in $k$ and agrees with the first at $u$.  Consequently
a minimum always occurs at $k=n$, with value
$1/n+u(n-1)/n^2$.  For $u=0$ the same value is $1/n$.
Taking the largest allowable $H$ makes
$u=\max\{n-M+1,0\}$ and proves
\eqref{eq:r78-strict-hierarchy}.  The construction uses a uniform
initial command row and one pure row for each action of $H$;
reports in $H$ select their pure row and the remaining reports
select the initial row.  The stated data and positivity properties
follow directly from the displayed primitive construction.

Finally, the Bellman residuals for this family are
$d_{ia}=\mathbf1_{\{a\ne i\}}$.  The matrix in
\eqref{eq:r77-reset-gap} retains only the first post-failure
penalty, so its minimized eigenvalue is obtained from
\eqref{eq:r78-cover-row-objective} by replacing $\lambda=1/n$
by $r/s=1/(2n)$.  The same row-purification argument applies.
For $k\le u$ the candidate is $1/k+(u-1)/(2n)$; for $k\ge u$
it is $u/(2n)+(1-u/(2n))/k$.  Again $k=n$ is a minimizer,
and subtracting $V_*=4+1/n$ gives
$D_M=u(n-1)/(2n^2)$, proving
\eqref{eq:r78-certificate-gap}.  This equality is specific to
the present family and is not a universal identity for $D_M$.

Continuity from Proposition~\ref{prop:r78-frontier-stability},
applied to the finitely many adjacent gaps, proves the neighborhood
assertion.  For an explicit perturbation put
$X=|e_1\rangle\langle f_1|+|f_1\rangle\langle e_1|$ and replace
only $E_{1,1}$ by $E_{1,1}+\epsilon X$, where
\[
 0<\epsilon<\frac{n-1}{8n^2}.
\]
The failure effect remains positive definite, and the corresponding
success effect $\tfrac12I-\epsilon X$ remains positive definite.
All other primitive data are unchanged.  Since
$G_1(e_1)=1$ and $G_1(f_1)=4$, this modified failure effect does
not commute with $G_1$.

These cost and report effects have a joint physical implementation.
Charge one on every failure.  Split the success effect $B_a$ into
\[
 F_{a,9}=(G_a-I)/8,\qquad F_{a,1}=B_a-F_{a,9},
\]
charging nine and one, respectively, and retaining only the common
success report at the control port.  Both effects are positive:
$0\preccurlyeq F_{a,9}\preccurlyeq3I/8$ and
$B_a\succeq(1/2-\epsilon)I$, with $\epsilon<1/32$.
Their sum is $B_a$, and the aggregate mean cost is
$I+8F_{a,9}=G_a$.  Each failure branch uses its stated
measure-and-prepare reset map, and each success subbranch can use
any terminal prepared state.  Thus the perturbation is an actual
instrument with positive, bounded paid costs and no additional
observation, rather than an unsupported independent assignment
of noncommuting cost and report operators.

The perturbation has zero expectation in every $\sigma_i$, so it
changes none of the reset transition probabilities or continuation
Bellman vectors for any chosen list of command rows.  Each such
vector satisfies $x_i\le4$, since reset costs are at most two and
each subsequent call succeeds with probability $1/2$.  Its initial
cost operator changes only by $\epsilon p^1_1x_1X$, of norm at
most $4\epsilon$.  Hence the optimized values satisfy
$|V_M(\epsilon)-V_M(0)|\le4\epsilon$ for every $M$.
Each adjacent gap is therefore at least
$(n-1)/n^2-8\epsilon>0$.  The common unrestricted upper memory
bound still gives $V_{n+1}(\epsilon)=V_*(\epsilon)$.  This proves
the claimed noncommuting realization without extending the
unperturbed closed formula to perturbed data.
\end{proof}

The family shows that the full frontier carries information absent
from its saturation threshold: each intermediate memory can have a
different, attained optimal value under a fixed success probability
and uniformly bounded primitive costs.  Formula
\eqref{eq:r78-cover-spectrum} retains the combinatorial difficulty
of the arbitrary incidence table, while giving a complete finite
answer rather than an efficient algorithm claim.  The reset
assumption and the declared observation port remain essential;
none of these statements completes the separate general stopping,
phase-type timing, or squared-sign minimax-risk problems.

\section{Exact optimal memory for reset experiments with tied actions}
\label{sec:r77-reset-memory}

The realization of a prescribed response and the implementation of an
optimal feedback rule are different problems.  We now give an exact
memory classification for a class of controlled experiments in which
optimal actions may be tied.  The class allows any finite command set,
any finite number of reset reports, arbitrary finite Hilbert dimension,
state-dependent call costs, and an unknown initial density matrix.
The additional hypothesis is strict positivity of each failure effect.
It supplies the reachability needed for a sharp converse; it is not a
uniqueness assumption on an optimal decision.

Use the reset instrument of Section~\ref{sec:reset-control}:
\[
 \mathcal I_{a,i}(\rho)=\tr(E_{a,i}\rho)\sigma_i,
 \qquad \sum_{i=1}^s E_{a,i}\preccurlyeq I,
 \qquad c_*I\preccurlyeq G_a\preccurlyeq c^*I.
\]
Here $a\in\mathcal A$, $1\le i\le s$, $c_*>0$, and success requires
immediate commitment.  There is no additional state measurement or
free reset.  Throughout this section assume
\begin{equation}\label{eq:r77-reset-positive}
 E_{a,i}\succ0\quad(a\in\mathcal A,\ 1\le i\le s),
 \qquad\text{there is no closed failure set.}
\end{equation}
Theorem~\ref{thm:reset-control} supplies the finite Bellman vector $v$,
the start operators $L_a=G_a+\sum_i v_iE_{a,i}$, and the value
\[
 V_* =\min_{p\in\Delta(\mathcal A)}
       \lambda_{\max}\Bigl(\sum_a p_aL_a\Bigr).
\]
Define the Bellman residuals and optimal action sets by
\[
 d_{ia}=\tr(G_a\sigma_i)+\sum_j\tr(E_{a,j}\sigma_i)v_j-v_i,
 \qquad \mathcal A_i^*=\{a:d_{ia}=0\}.
\]
Every $d_{ia}$ is nonnegative and every $\mathcal A_i^*$ is nonempty.
Put
\[
 \mathcal P_* =\left\{p\in\Delta(\mathcal A):
       \lambda_{\max}\Bigl(\sum_a p_aL_a\Bigr)=V_*\right\},
 \qquad I(p)=\{i:\supp p\subseteq\mathcal A_i^*\}.
\]
For $J\subseteq\{1,\ldots,s\}$ let
\begin{equation}\label{eq:r77-transversal}
 t(J)=\min\{|H|:H\subseteq\mathcal A,
               \ H\cap\mathcal A_i^*\ne\varnothing\ (i\in J)\},
 \qquad t(\varnothing)=0.
\end{equation}
This is an ordinary transversal number of the displayed finite set
system.  It counts distinct continuing command rows, rather than
physical states or the support size of a least-favorable preparation.

\begin{theorem}[Exact optimal-memory formula with ties]
\label{thm:r77-reset-cover}
Under \eqref{eq:r77-reset-positive}, the least number of live labels of
a fixed-program, prescribed-pure-start autonomous controller attaining
the all-initial-state optimum $V_*$ is
\begin{equation}\label{eq:r77-reset-memory}
 m_* =1+\min_{p\in\mathcal P_*}
                  t\bigl(\{1,\ldots,s\}\setminus I(p)\bigr).
\end{equation}
Every minimum is realized by one optimal initial command row and pure
continuing command rows.  The post-report transition may be chosen as a
deterministic function of the current reset report alone.  Arbitrary
stochastic transitions, dependence on the previous label and command,
and mixed continuing rows cannot lower \eqref{eq:r77-reset-memory}.

Event stopping has the same total-label minimum.  State stopping has
total-label minimum $m_*+1$.  For algebraic primitive data the value,
optimal label count, and one realizing controller can be determined
by a finite exact procedure; polynomial running time is not asserted.
\end{theorem}
\begin{proof}
Choose an attaining dual density matrix $\rho_*$ from
Theorem~\ref{thm:reset-control}.  Consider any robust optimal controller,
write its fixed command rows as $p^1,\ldots,p^M$, and index its pure
initial label by $1$.  Its cost from every $\rho$ is at least
$\tr(\rho\sum_a p^1_aL_a)$.  Consequently
$\lambda_{\max}(\sum_a p^1_aL_a)\le V_*$, and the reverse inequality
holds by the definition of $V_*$.  Thus $p^1\in\mathcal P_*$.

The Bellman residual identity of Proposition~\ref{prop:reset-rigidity}
holds at $\rho_*$.  Its initial term is at least $V_*$, whereas the
whole cost is at most $V_*$.  Every nonnegative residual on a reached
continuation therefore vanishes.  For every reset report $i$, its
probability on the first call is
\[
 \sum_a p^1_a\tr(E_{a,i}\rho_*)>0.
\]
This strict inequality follows from positive definiteness and holds
even if $\rho_*$ is singular or $p^1$ has proper support.  After this
report at least one label $k$ is reached with positive probability.
The next command has row $p^k$, so
$\sum_a p^k_a d_{ia}=0$.  Nonnegativity implies
$\supp p^k\subseteq\mathcal A_i^*$.

If $i\notin I(p^1)$, that label cannot be $1$.  For each other label
which serves one or more such states, choose any action in the support
of its row.  It belongs to every optimal action set served by that row.
These at most $M-1$ chosen actions meet all
$\mathcal A_i^*$ with $i\notin I(p^1)$.  This proves the lower bound
in \eqref{eq:r77-reset-memory}.  It uses only first-failure events;
no reachability conclusion is inferred from an unrelated path law.

Conversely, fix $p\in\mathcal P_*$ and a smallest transversal $H$ of
the remaining states.  Use one initial label with command row $p$ and
one pure command label for each $a\in H$.  After a failure report $i$,
return to the initial label when $i\in I(p)$; otherwise select a fixed
$a\in H\cap\mathcal A_i^*$.  Every continuing row is then Bellman
optimal for its actual reset state.  Bellman's equality bounds every
finite-prefix cost from $\sigma_i$ by $v_i$.  The positive cost floor
implies finite expected duration, and the bounded remaining potential
vanishes.  Hence the continuation cost is exactly $v_i$.  From an
arbitrary initial state the full cost is
$\tr(\rho\sum_a p_aL_a)\le V_*$.  This is a proper robust optimum.
There are finitely many supports of $p$, so the integer minimum in
\eqref{eq:r77-reset-memory} is attained.  Reusing a live label for the
single terminal decoder is permitted exactly in the event-stopped
interface, giving the claimed terminal counts.

For the effective assertion enumerate deterministic command selectors
on the $s$ reset states.  A proper selector has cost
$(I-Q_\alpha)^{-1}c_\alpha$; properness and this vector can be checked
by finite algebraic operations and spectral or graph tests.  Some
proper selector is simultaneously optimal by
Theorem~\ref{thm:reset-control}, so the enumeration determines $v$.
The spectral minimization defining $V_*$ is a compact semialgebraic
optimization.  For each nonempty support $T\subseteq\mathcal A$,
the assertion that some $p\in\mathcal P_*$ has exactly that support
is a first-order real-algebraic condition.  For example, impose
$p_a>0$ on $T$, $p_a=0$ off $T$, $\sum_a p_a=1$, and
$\sum_a p_aL_a\preccurlyeq V_*I$; positivity of a Hermitian matrix
is expressed by its principal minors after separating real and
imaginary parts.  Exact elimination decides this condition and
supplies a sample when it is true.  Finite transversal enumeration
then gives the minimum and the stated implementation.  The procedure
does not assume that a numerical optimizer has certified a tie.
\end{proof}

The positive-effect hypothesis can be weakened to the existence of one
attaining $\rho_*$ for which $\tr(E_{a,i}\rho_*)>0$ for all $a,i$;
the same proof applies.  The stated version is a direct primitive
condition, without an unknown optimal-trajectory assumption.  When
some reports can be inaccessible, the formula for all $s$ reset
states need not be necessary: an unreachable state may carry an
arbitrary optimal action set.  In that case the lower certificates
of Proposition~\ref{prop:reset-rigidity} retain their original scope.

\subsection{A finite gap certificate for every insufficient memory}

For $M\ge1$ and rows $p^1,\ldots,p^M\in\Delta(\mathcal A)$ put
\[
 r_i(p^1,\ldots,p^M)=\min_{1\le k\le M}\sum_a p^k_a d_{ia},
 \qquad
 K_M(p^1,\ldots,p^M)=\sum_a p^1_aL_a
       +\sum_{a,i}p^1_a r_i(p^1,\ldots,p^M)E_{a,i},
\]
and define the finite variational quantity
\begin{equation}\label{eq:r77-reset-gap}
 D_M=\min_{(p^1,\ldots,p^M)\in\Delta(\mathcal A)^M}
                  \lambda_{\max}(K_M(p^1,\ldots,p^M))-V_*.
\end{equation}

\begin{theorem}[Attained certificate of a strict optimality gap]
\label{thm:r77-reset-gap}
Under \eqref{eq:r77-reset-positive}, every robustly proper controller
with at most $M$ live labels has worst-initial-state expected cost at
least $V_*+D_M$.  The minimum defining $D_M$ is attained, and
\[
 D_M>0\quad\Longleftrightarrow\quad M<m_*.
\]
For algebraic primitive data, $D_M$ is computable by exact real
algebraic optimization.  It is a lower certificate, not a claimed
formula for the entire memory-constrained cost frontier.
\end{theorem}
\begin{proof}
For the fixed rows of a controller, condition on the first command and
its failure report $i$.  Whichever label the controller selects, the
expected Bellman residual on its next call is at least $r_i$.  All
later residuals are nonnegative.  The residual identity, now at any
initial $\rho$, gives cost at least $\tr(\rho K_M)$, and maximization
over density matrices gives $\lambda_{\max}(K_M)$.  A controller of
infinite worst-state cost satisfies the lower bound trivially.

The function in \eqref{eq:r77-reset-gap} is continuous on a compact
product of simplexes.  Moreover $K_M\succeq\sum_a p^1_aL_a$, so
$D_M\ge0$.  If $D_M=0$, take minimizing rows.  Then $p^1\in\mathcal
P_*$.  Evaluate $K_M$ at an attaining dual $\rho_*$.  Its first term
is at least $V_*$ and each coefficient
$\sum_a p^1_a\tr(E_{a,i}\rho_*)$ is strictly positive.  Since its
value is at most $\lambda_{\max}(K_M)=V_*$, every $r_i$ is zero.
Thus every reset state has at least one row supported on its optimal
action set.  The construction in Theorem~\ref{thm:r77-reset-cover}
implements an optimal controller with these rows, hence $M\ge m_*$.
Conversely an optimal controller supplied there has $r_i=0$ for all
$i$, so $D_M=0$ whenever $M\ge m_*$.  Compactness makes the remaining
cases strictly positive.  Finite minima, products and Hermitian
eigenvalue inequalities are semialgebraic, proving the last assertion.
\end{proof}

\subsection{Every finite optimal-action cover has a physical realization}

The preceding set system is not merely an artifact of a bound.
Arbitrary finite covering problems occur in reset instruments with
commuting cost and effect operators
with strictly positive failure effects and constant success probability.

\begin{theorem}[Universality of the optimal-memory covering problem]
\label{thm:r77-reset-universality}
Let $q\ge2$ and let $A_1,\ldots,A_s$ be nonempty proper subsets of
$\{1,\ldots,q\}$.  Let $t$ be their transversal number.  For every
fixed rational $0<r<1$ there is an explicit rational reset
instrument with commuting cost and effect operators in dimension $s+q$, with $q$ commands, $s$ failure reports,
success probability $1-r$ at every call, and costs in $[1,4]$, such
that
\[
 V_*=3+\frac{r}{1-r}+\frac1q,
 \qquad m_*=t+1.
\]
Its unique optimal initial command row is uniform.  The minimum is
over all stochastic autonomous controllers with a prescribed pure
initial label, and includes arbitrary randomized post-report updates.
Thus determining optimal memory is NP-hard even on this class, while
the unrestricted value and an optimal command distribution are explicit.
\end{theorem}
\begin{proof}
Take an orthonormal basis $e_1,\ldots,e_s,f_1,\ldots,f_q$ and set
\[
 \sigma_i=|e_i\rangle\langle e_i|,
 \qquad E_{a,i}=\frac r s I.
\]
The failure map is $\mathcal I_{a,i}(\rho)=(r/s)\tr(\rho)\sigma_i$.
For instance its Kraus operators are
$\sqrt{r/s}\,|e_i\rangle\langle b|$ for all displayed basis vectors
$b$.  Add any success measure-and-prepare branch with effect $(1-r)I$.
Their sum is trace preserving, so these are actual quantum instruments.
They commute at the level of the relevant effects and all post-failure
states are classical basis states.

Specify the positive diagonal cost operator $G_a$ by
\[
 \langle e_i,G_ae_i\rangle=1+\mathbf1_{\{a\notin A_i\}},
 \qquad
 \langle f_j,G_af_j\rangle=3+\mathbf1_{\{a=j\}}.
\]
These means are physically realizable without extra controller
information.  In each call measure the displayed diagonal basis,
charge its indicated cost, then use fresh coins to choose success or
one of the reset reports and prepare the corresponding state.  The
intermediate basis outcome and realized cost are not an additional
observation at the control port.  Averaging them gives exactly the
declared instrument and mean cost; no coherence-dependent hidden
assumption is needed.

All commands have $Q_{iaj}=r/s$.  Since every $A_i$ is nonempty,
$v_i=(1-r)^{-1}$ solves the Bellman equations and is their finite
solution.  Its optimal action set at $i$ is exactly $A_i$.
The start operator is $L_a=G_a+r(1-r)^{-1}I$.  For any initial
command row $p$, its eigenvalues on the $e_i$ coordinates are
\[
 \frac1{1-r}+p(A_i^c)\le2+\frac r{1-r},
\]
whereas its eigenvalues on the $f_j$ coordinates are
$3+r(1-r)^{-1}+p_j$.  The latter strictly dominate, hence
\[
 \lambda_{\max}\Bigl(\sum_a p_aL_a\Bigr)
       =3+\frac r{1-r}+\max_j p_j.
\]
Its unique minimum is the uniform row, with the displayed value.
Its full support lies in none of the proper sets $A_i$, so
Theorem~\ref{thm:r77-reset-cover} gives $m_*=1+t$ and constructs an
attaining controller.

The matrices have dimension and rational encoding length polynomial
in the incidence table of the set system; one may fix $r=1/2$.
Minimum transversal is the classical covering problem
\cite{Karp72}.  Restriction to proper nonempty subsets does not remove
its hardness: adjoining an unused action makes every nonempty set
proper without changing its transversal number.  An oracle for the
optimal label number in the displayed physical family would therefore
return the transversal number by subtracting one.  This proves the
complexity assertion by an explicit reduction, not by a general
complexity claim about POMDPs.
\end{proof}

\begin{example}[Why positive reachability cannot be dropped]
\label{ex:r77-reset-exposure}
On $\mathbb C^2$ take one reset state $\sigma=|0\rangle\langle0|$,
two commands $a,b$, and
\[
 E_a=E_b=\tfrac12|0\rangle\langle0|,
 \qquad G_a=\operatorname{diag}(1,6),\qquad
 G_b=\operatorname{diag}(2,5).
\]
There is no closed failure set.  The Bellman value is $v=2$ and its
unique optimal continuation command is $a$.  The start operators
are $L_a=\operatorname{diag}(2,6)$ and
$L_b=\operatorname{diag}(3,5)$, so $V_*=5$ and the unique minimizing
initial row is the pure command $b$.  Applying
\eqref{eq:r77-reset-memory} without positive reachability would give
two labels.  Nevertheless the constant $b$ controller has cost
operator $\operatorname{diag}(4,5)$: from $\sigma$ it pays $2$ per
call with failure probability $1/2$, and from $|1\rangle$ it pays
$5$ and succeeds immediately.  It is a one-label robust optimum.
The least-favorable preparation $|1\rangle\langle1|$ never sees the
reset report.  Thus a locally nonoptimal continuation can be hidden
below the robust optimum when the exposing preparation cannot reach
it.  This verifies the need for the reachability hypothesis by exact
primitive data.
\end{example}

For example, the three optimal sets $\{1,2\}$, $\{2,3\}$ and
$\{1,3\}$ give $m_*=3$, although every report state has two optimal
commands.  Pairwise uniqueness of Bayes actions plays no role.  More
generally the count can grow with an arbitrary finite covering
instance while the success probability remains fixed and every
failure effect remains strictly positive.

Bellman verification and shortest-path optimization are classical
\cite{BertsekasTsitsiklis91}; search over finite stochastic controllers
also has an extensive history \cite{Meuleau99,PoupartBoutilier03}.
The assertions proved here concern the exact counted-row invariant,
its initializer interaction, the strict finite gap certificate and
its realization inside the stated reset experiment.  They do not
assert a general classification of partially observed control or
replace the distinct squared-sign minimax-risk problem.

\section{Reversible response groups and common controlled memory}\label{sec:r72-reversible}
The fixed-phase results concern one transition matrix. Several commands require one family of matrices on the same retained labels, and checking every command in isolation is insufficient. We give an intrinsic finite-realization criterion for responses whose minimal linear dynamics is reversible and compact, including any finite family of unitary commands on a prepared finite-dimensional quantum system. We then obtain a matching arithmetic memory law for a pair of noncommuting involutions. This is a theorem about exact controlled response realization, not an assertion that a simulator's private label is an observed sufficient statistic for arbitrary feedback optimization.

Positive realization through invariant polyhedral sets is classical \cite{Heller65,MonrasWinter14,ThonJaeger15}. The argument below uses a section of the realizing simplex before passing to the observable quotient; auxiliary hidden-state responses need not belong to the minimal response space. Compact reversibility then turns containment into equality. The resulting finite-group obstruction concerns all command words, rather than a single spectral restriction.

\subsection{The minimal response space}
Let $A$ and $C$ be nonempty finite alphabets of commands and terminal outcomes. A supplied response is a function $b:A^*\to\Delta(C)$. Its residual after a word $w$ is $b^w(z)=b(wz)$, including the empty word $\varnothing$. Put
\[
 V_b=\operatorname{span}_{\R}\{b^w:w\in A^*\},
 \qquad (L_af)(z)=f(az).
\]
Assume $V_b$ is finite dimensional. It is invariant under each $L_a$, since $L_ab^w=b^{wa}$. Call the response \emph{compactly reversible} if the restrictions $L_a|_{V_b}$ are invertible and the closure $G_b$ of their generated group is a compact subgroup of $\operatorname{GL}(V_b)$. This concerns the observable quotient; it does not require the full physical command group to be finite or fully observable.

An exact stochastic realization with $M$ live labels has row-stochastic matrices $Q_a$, one row-stochastic terminal decoder matrix $D$, and the prescribed initial row $e_1^{\mathsf T}$, such that
\begin{equation}\label{eq:r72-response}
 b(a_1\cdots a_n)=e_1^{\mathsf T}Q_{a_1}\cdots Q_{a_n}D.
\end{equation}
Randomized terminal readout is allowed even when all $Q_a$ are deterministic. The target input word is not encoded into an uncounted clock or persistent seed.

\begin{theorem}[Finite realization of compactly reversible responses]\label{thm:r72-reversible}
For a compactly reversible response, the following are equivalent:

\textup{(i)} An exact finite stochastic realization exists.

\textup{(ii)} The residual set $\{b^w:w\in A^*\}$ is finite.

\textup{(iii)} The observable group $G_b$ is finite.

\textup{(iv)} An exact realization with deterministic transitions and a possibly randomized decoder exists.

The least number of deterministic-transition labels is the residual cardinality. This is not in general the least stochastic label count. If an $M$-label stochastic realization exists, then
\[
 |G_b|\le(2^M-1)!.
\]
This last bound is only a finite combinatorial certificate, not a sharp memory law.
\end{theorem}
\begin{proof}
Suppose \eqref{eq:r72-response} holds, and write $r_w=e_1^{\mathsf T}Q_w$. Let $R=\operatorname{span}_{\R}\{r_w\}\subseteq\R^M$. The assignment $\pi(r_w)=b^w$ extends to a well-defined surjective linear map $\pi:R\to V_b$: a relation $\sum_wc_wr_w=0$ implies, after right multiplication by $Q_zD$, the identical relation among $b^w(z)$ for every $z$. Moreover
\[
 \pi(rQ_a)=L_a\pi(r)\qquad(r\in R).
\]
The simplex section $K=R\cap\Delta_M$ is a compact polytope, invariant under right multiplication by every $Q_a$. Its image $P=\pi(K)$ is a compact polytope, contains every residual, and spans $V_b$. It satisfies $L_aP\subseteq P$. It is important that we intersect the simplex with $R$ first: projecting arbitrary hidden-state responses directly into $V_b$ would not be justified.

For an element $T$ of a compact matrix group, positive powers have a subsequence tending to the identity with exponents tending to infinity. One elementary proof takes a convergent subsequence of powers with successively growing index gaps; quotients of two close powers approach the identity. A finite-order element is covered by its period multiples. Thus positive powers also approach $T^{-1}$. Applying this to $L_a$ and using closedness of $P$ gives $L_a^{-1}P\subseteq P$, and hence $L_aP=P$. All elements of $G_b$ preserve $P$. They permute its finitely many vertices, and this permutation action is faithful because those vertices span $V_b$. Consequently $G_b$ is finite. This proves (i)$\Rightarrow$(iii).

For the displayed bound, a vertex of $K$ is determined by its positive support. If two distinct vertices had the same support, their difference could be added and subtracted in sufficiently small multiples while preserving $R$, unit mass and nonnegativity, contradicting extremality. Thus $K$ has at most $2^M-1$ vertices. A linear image of a polytope has no more vertices than the original polytope, so the faithful vertex action proves the bound.

If $G_b$ is finite, the residual orbit is finite. Conversely, if the residual set is finite, each invertible $L_a$ maps it injectively into itself, hence permutes it. Since it spans $V_b$, the generated group and its closure are finite. This proves (ii)$\Leftrightarrow$(iii). Under (ii), use the residuals as labels, send $b^w$ to $b^{wa}$ on input $a$, and use decoder $b^w(\varnothing)$. This realizes $b$ with deterministic transitions. If any such deterministic-transition machine sends $u$ and $v$ to the same label, every suffix yields the same decoder law, so $b^u=b^v$. This proves its minimal count. Finally (iv)$\Rightarrow$(i) is immediate. No assertion that stochastic transitions can always be determinized without increasing size was used.
\end{proof}

\begin{corollary}[Prepared unitary commands]\label{cor:r72-unitary}
Let $U_a$ be finitely many unitary matrices on $\mathbb C^d$, let $\rho$ be a prepared density matrix, and let $\{E_r:r\in C\}$ be a terminal positive-operator-valued measurement. For
\[
 b(w,r)=\tr(E_r\Phi_w(\rho)),\qquad \Phi_a(X)=U_aXU_a^*,
\]
exact finite stochastic realization for all command words is possible if and only if its observable response group is finite. The group is computed after quotienting only those directions invisible to every future commanded terminal test. Its dimension is at most $d^2$.
\end{corollary}
\begin{proof}
Let $W=\operatorname{span}_{\R}\{\Phi_w(\rho)\}$ in the Hermitian matrices. Each conjugation is a Hilbert--Schmidt isometry. Its forward invariant subspace $W$ is also backward invariant, by finite-dimensional invertibility. Let
\[
 N=\{X\in W:\tr(E_r\Phi_zX)=0\text{ for all }z,r\}.
\]
It is invariant under every $\Phi_a$ and hence under its inverse. The orthogonal complement of $N$ in $W$ is invariant as well. Therefore the induced transformations on $W/N$ are orthogonal for the quotient inner product. The map $X\mapsto(z,r\mapsto\tr(E_r\Phi_zX))$ identifies $W/N$ with $V_b$ and intertwines them with $L_a$. Their generated group has compact closure in the orthogonal group. Apply Theorem~\ref{thm:r72-reversible}. A constant measurement can make this induced group trivial even when the physical unitary group is infinite, so finiteness of the full physical group is not a necessary condition.
\end{proof}

\subsection{A matching memory law for two noncommuting commands}
Fix $x\in\R$ with $\sin(2\pi x)\ne0$. On a protected qubit take
\[
 U_0=\begin{pmatrix}1&0\\0&-1\end{pmatrix},\qquad
 U_1=\begin{pmatrix}\cos(\pi x)&\sin(\pi x)\\
                     \sin(\pi x)&-\cos(\pi x)\end{pmatrix},
 \quad \rho=E=|0\rangle\langle0|.
\]
Both commands are involutions. For known $0\le\alpha\le1$, their channels are
\[
 \Phi_{a,\alpha}(X)=\alpha U_aXU_a^*+(1-\alpha)\tr(X)I/2.
\]
Let $M_x^{\rm ctrl}(\alpha)$ be the least stochastic live memory of the whole binary response $w\mapsto\tr(E\Phi_{w,\alpha}\rho)$, using the same matrices and decoder for every word. In the real Bloch plane, with coordinates $(z,x)$, these commands are $\alpha S_0$ and $\alpha S_1$, where $S_0$ reflects the angle $t$ to $-t$ and $S_1$ reflects it to $2\pi x-t$. Their product is $\alpha^2R_{2\pi x}$.

\begin{theorem}[All-word arithmetic law for an involutive command pair]\label{thm:r72-control-envelope}
For $0<\epsilon<1$ and $\alpha=1-\epsilon$, the arithmetic envelope of Section~\ref{sec:r71-arithmetic} gives
\begin{equation}\label{eq:r72-control-envelope}
 \tfrac14\mathfrak A_x(\epsilon)
 \le M_x^{\rm ctrl}(1-\epsilon)
 \le(1+4\pi^2)\mathfrak A_x(\epsilon).
\end{equation}
The upper is one explicit finite stochastic machine valid for every mixed command word, not separately chosen unary machines. The lower ranges over arbitrary common stochastic realizations.

For irrational $x$, the lower and upper logarithmic growth exponents as $\epsilon\downarrow0$ are respectively
\[
 \frac1{\mu(x)+1},\qquad\frac{\mu(x)-1}{2\mu(x)-1},
\]
with values $0,1/2$ at $\mu=\infty$. A uniform cubic-root law holds precisely for badly approximable $x$; otherwise the normalized count has lower limit zero and upper limit infinity.

At $\alpha=1$, the exact count is the reduced denominator $q$ if $x$ is rational, and is infinite otherwise. At $\alpha=0$, it is two. Every strictly damped pair has a finite exact common realization. These statements concern controlled response memory, not a general multi-command boundary-optimal policy theorem.
\end{theorem}
\begin{proof}
Any common realization restricted to the words $(01)^n$ realizes
\[
 \tfrac12[1+\alpha^{2n}\cos(2\pi nx)]
\]
with transition matrix $Q_0Q_1$. Its size is therefore at least
$m_x(1-\alpha^2)\ge\mathfrak A_x(\epsilon(2-\epsilon))/2$
by Theorem~\ref{thm:r71-envelope}. For $1\le c\le2$, the defining maximum of the envelope obeys $\mathfrak A_x(c\epsilon)\ge\mathfrak A_x(\epsilon)/c$ whenever both arguments are in its domain. This follows for each denominator before taking its infimum. Taking $c=2-\epsilon$ proves the lower in \eqref{eq:r72-control-envelope}.

For the upper, the rational multiple-polygon construction in the proof of Theorem~\ref{thm:r71-envelope} gives a regular $l$-gon $P$ on the unit circle, with vertex $v_0=(1,0)$, such that
\[
 l\le(1+4\pi^2)\mathfrak A_x(\epsilon),\qquad
 \alpha R_{2\pi x}v_j\in P\quad\text{for every vertex }v_j.
\]
The first reflection $S_0$ permutes this regular polygon and zero belongs to it. Consequently $\alpha S_0P\subseteq P$, while $S_1=R_{2\pi x}S_0$ gives $\alpha S_1P\subseteq P$. Choose barycentric rows $Q_a(j,k)$ of $\alpha S_av_j$ in the same vertices. Let $d_j=(1+(v_j)_1)/2$. Starting at vertex zero, induction on the actual command word shows that the expected retained vertex equals the protected Bloch vector. Pairing with the first coordinate proves every response identity. The start is one of the counted vertices; there is no extra initialization mixture.

The two constants in \eqref{eq:r72-control-envelope} transfer both exponents and both uniform-law statements from Theorems~\ref{thm:r71-exponents} and \ref{thm:r71-cubic-rigidity}. At zero damping, a rational phase of reduced denominator $q\ge3$ has a common deterministic $q$-gon realization. The word-pair transition has a primitive $q$th-root mode and needs at least $q$ states by the unit-circle spectral obstruction. For irrational $x$ that same obstruction rules out every finite size. At $\alpha=0$ the response is one at the empty word and $1/2$ after every nonempty word, realized by a first label followed by an absorbing fair label; one label cannot change its decoder. This proves all endpoints.
\end{proof}

A fully algebraic example uses $U_1=5^{-1/2}\left(\begin{smallmatrix}2&1\\1&-2\end{smallmatrix}\right)$. Then $e^{2\pi ix}=(3+4i)/5$ is not a root of unity, since its sum with its conjugate is the rational noninteger $6/5$. Each undamped command in isolation has period at most two; their common all-word response has no finite stochastic realization. The arithmetic exponent of this particular angle is not determined here. The example isolates the failure of commandwise testing without inferring a Diophantine class from finitely many coefficients.

If an independent herald with fixed probability $p$ is placed before each command operation, and only failures apply the supplied command, every common realization above preserves the terminal response under every external command schedule and the first-herald cost $1/p$. The commands are actual inputs to that same retained-state transducer. This does not make its private label observable to the caller, allow uncharged midstream quantum measurements, or prove optimality for a controller free to choose an easier command sequence. The unary optimal squared-task interpretation below continues to use its separately proved physical embedding.

\section{Convex specifications and optimal uniform approximation}
\label{sec:r76-convex-specifications}

The exact normal form has a robust extension. The target need not be an
exactly specified probability table, and the stochastic controller need
not satisfy the group relations, even on its observable states. What is
required is a constraint on every word representing a given group element.
The quantifier over all words is essential.

Fix a finite command alphabet $A$, a finite group $G$, and a surjective
homomorphism $\gamma:A^*\to G$. For every $g\in G$ let $K_g$ be a closed
convex subset of the finite probability simplex $\Delta(C)$. Empty sets
are allowed. An $M$-label pure-start controller consists of a prescribed
point row $e_o$, stochastic command matrices $Q_a$, and a row-stochastic
decoder $D_0$. It satisfies the specification if
\begin{equation}\label{eq:r76-specification}
 e_oQ_wD_0\in K_{\gamma(w)}\qquad(w\in A^*).
\end{equation}
For a mixed start, replace $e_o$ by an arbitrary probability row and do
not charge an additional label for preparing that row.

\begin{theorem}[Finite-group convex-specification normal form]
\label{thm:r76-convex-normal-form}
The specification \eqref{eq:r76-specification} has a pure-start
realization on at most $M$ labels if and only if there are a finite
$G$-set $X$, a probability row $\lambda$ on $X$, and a row-stochastic
decoder $D$ such that
\begin{equation}\label{eq:r76-convex-core}
 \lambda P_gD\in K_g\quad(g\in G),\qquad
 |X|+\mathbf1_{\{\lambda\text{ is not a point mass}\}}\le M.
\end{equation}
For mixed starts the corresponding criterion is $|X|\le M$.
The action on $X$ is permitted to have a kernel. In particular no
faithfulness hypothesis is imposed on an approximating action.
\end{theorem}
\begin{proof}
We give the projection argument explicitly because its use of convexity
replaces the exact observable intertwining in
Theorem~\ref{thm:r73-normal-form}.
Let $\mathcal L=\gamma^{-1}(e)$. Select finitely many identity words
$\ell_1,\ldots,\ell_s$ whose positive matrix entries cover every
directed edge occurring in any $Q_\ell$, $\ell\in\mathcal L$, and put
\[
 T=(I+Q_{\ell_1}+\cdots+Q_{\ell_s})/(s+1).
\]
At most $M^2$ words suffice. Each closed communicating class $C_i$ of
$T$ is closed under every identity-word matrix. The positive diagonal
and Lemma~\ref{lem:r73-lazy-convergence} give a stochastic idempotent
$E=\lim_nT^n$. Its image simplex has disjointly supported vertices
$\mu_1,\ldots,\mu_k$, the stationary rows of the closed classes.
Every identity word preserves those classes, and consequently
\begin{equation}\label{eq:r76-sandwich-identity}
 EQ_\ell E=E\qquad(\ell\in\mathcal L).
\end{equation}
For a word $u$ write $F_u=EQ_uE$ on this image simplex. Choose an
inverse word $\bar u$ with $\gamma(\bar u)=\gamma(u)^{-1}$.
Since $T^n$ is a convex combination of identity-word matrices,
\eqref{eq:r76-sandwich-identity} applied to $u\ell\bar u$ and
$\bar u\ell u$ gives
\[
 F_uF_{\bar u}=E=F_{\bar u}F_u.
\]
Thus $F_u$ is an affine automorphism of the image simplex and permutes
its vertices. If $\gamma(u)=\gamma(v)$, multiplication by the common
inverse $F_{\bar u}$ shows $F_u=F_v$. Replacing the middle $E$ by
$T^n$ in $EQ_uEQ_vE$ also gives $F_uF_v=F_{uv}$, since all words
$u\ell v$ have the same group image. This defines a $G$-action $P_g$
on $X=\{1,\ldots,k\}$.

Write $e_oE=\sum_i\lambda_i\mu_i$ and put $D(i,\cdot)=\mu_iD_0$.
For any representative $u$ of $g$,
\[
 \lambda P_gD=e_oEQ_uED_0
             =\lim_{n\to\infty}e_oT^nQ_uT^nD_0.
\]
The last expression before passage to the limit is a convex combination
of $e_oQ_{\ell u\ell'}D_0$ with $\ell,\ell'\in\mathcal L$.
All these probability rows belong to $K_g$. Closedness and convexity
therefore give $\lambda P_gD\in K_g$. This argument includes $u$
empty and makes no claim that the original controller's word laws
factor through $G$.

If the initial label lies in a closed class $C_i$, then $e_oE=\mu_i$
and $\lambda$ is a point mass. Otherwise the $k$ disjoint nonempty
closed classes exclude at least the initial label, so $k+1\le M$.
This proves the pure-start charge. For a mixed row $\eta$, use
$\eta E$ instead and retain $k\le M$.

Conversely, a point preparation on $X$ needs only the permutation
transitions. For a mixed $\lambda$, a further label $\star$ has decoder
$\lambda D$ and, on the first actual command $a$, moves into $X$ with
row $\lambda P_{\gamma(a)}$. Subsequent transitions are permutations.
Its empty and nonempty word laws are respectively $\lambda D$ and
$\lambda P_{\gamma(w)}D$. This gives the required realization and
finishes both start conventions.
\end{proof}

For an exact probability table $b$, let $\kappa_G(b)$ be the least
$|X|$ among factorizations $b(g)=\lambda P_gD$. This is its mixed-start
minimum by the preceding theorem. The notation $m(b)$ retains the
counted pure-start convention; its exact formula and the residual
cardinality $N$ are given in Theorem~\ref{thm:r73-memory-formula}.

For a target $b:G\to\Delta(C)$ write
$\operatorname{TV}(p,q)=\tfrac12\sum_c|p_c-q_c|$ and define
\[
 \mathcal E_M^{\mathrm p}(b)=
 \inf_{\substack{\text{pure-start controllers on at most }M\text{ labels}}}
 \ \sup_{w\in A^*}\operatorname{TV}
           (b(\gamma(w)),e_oQ_wD_0).
\]
Use $\mathcal E_M^{\mathrm m}$ for the mixed-start convention.
For a finite $G$-set $X$ put
\[
 d_X(b)=\min_{\lambda,D}\max_{g\in G}
          \operatorname{TV}(b(g),\lambda P_gD),\qquad
 d_X^{\mathrm{pt}}(b)=\min_{x,D}\max_{g\in G}
          \operatorname{TV}(b(g),e_xP_gD).
\]
The decoder ranges over stochastic rows; compactness makes both minima
attained. Minima over no candidate actions are interpreted as $+\infty$.

\begin{theorem}[Attained finite optimization for infinite-word error]
\label{thm:r76-minimax-normal-form}
For every target $b$, finite $G$, and integer $M\ge1$,
\begin{align}
 \mathcal E_M^{\mathrm m}(b)&=\min_{|X|\le M}d_X(b),\label{eq:r76-mixed-error}\\
 \mathcal E_M^{\mathrm p}(b)&=
 \min\left\{\min_{|X|\le M}d_X^{\mathrm{pt}}(b),
             \min_{|X|\le M-1}d_X(b)\right\}.
 \label{eq:r76-pure-error}
\end{align}
One representative of each finite $G$-set is sufficient, so these are
finite optimizations and their minima are attained. Their values are
$1$-Lipschitz functions of $b$ in $\max_g\operatorname{TV}$.
If $\mathcal E_M^{\mathrm p}(b)>0$, a sufficiently small perturbation of
$b$ still cannot be exactly realized on $M$ pure-start labels; the same
assertion holds for mixed starts. When the target table is algebraic,
the optimum and a realizing finite controller are computable by real
quantifier elimination. This is a decidability assertion, not a
polynomial-time bound or a claim that quantifier elimination was executed.
\end{theorem}
\begin{proof}
Apply Theorem~\ref{thm:r76-convex-normal-form} to the closed convex
balls $K_g=\{p:\operatorname{TV}(p,b(g))\le\delta\}$.
It sends every controller of error at most $\delta$ to an admissible
action without increasing the label charge or the error. Conversely
every displayed action supplies such a controller. The pure constraint
splits precisely into a point preparation with $|X|\le M$ or an
arbitrary preparation with $|X|\le M-1$.
There are finitely many homomorphisms from the fixed finite group into
each symmetric group of degree at most $M$. For each action the allowed
preparations and decoders form compact sets. Hence the finite minima
are attained, including limits in which a mixed preparation becomes a
point mass. The triangle inequality proves the Lipschitz bound and
the perturbation assertion.

For the algorithmic statement, enumerate these actions and introduce
variables $\lambda_x,D_{xc},\delta$ and $u_{gc}\ge0$ with
\[
 -u_{gc}\le b(g,c)-(\lambda P_gD)_c\le u_{gc},\qquad
 \sum_cu_{gc}\le2\delta.
\]
Together with the stochastic constraints these are polynomial
inequalities over the real algebraic numbers. The point-start branch
fixes $\lambda=e_x$ and enumerates $x$. The compact attained minimum
is the least real point of the resulting semialgebraic set of feasible
$\delta$'s. Real quantifier elimination and algebraic sample-point
construction, used with the same scope as
Theorem~\ref{thm:r73-algebraic-synthesis}, compute the value and tables.
\end{proof}

The finite optimization concerns synthesis against a supplied finite-group
target. It does not solve the cutpoint, equivalence, or isolation problem
for a supplied arbitrary probabilistic automaton. Nor does it replace
the specification for all words by a check of one representative per group
element on the original stochastic controller. Those restrictions would
change the problem.

\begin{theorem}[Robust multiplicity and initialization certificates]
\label{thm:r76-robust-multiplicity}
Use the stacked Fourier matrix $A_\rho(b)$ of
\eqref{eq:r75-fourier-block}, with singular values in decreasing order,
and write $m_\rho(X)$ for the multiplicity of $\rho$ in
$\mathbb C[X]$. If an action $X$ approximates $b$ to total-variation
error at most $\delta$, then
\begin{equation}\label{eq:r76-fourier-robust}
 \sigma_{m_\rho(X)+1}(A_\rho(b))\le\sqrt2\,\delta
                       \qquad(\rho\in\widehat G).
\end{equation}
A singular value past the number of columns means zero. For a binary
response $b=(\beta,1-\beta)$ the sharper equivalent scalar certificate is
\begin{equation}\label{eq:r76-binary-robust}
 \sigma_{m_\rho(X)+1}(\widehat\beta(\rho))\le\delta.
\end{equation}
In particular let $r_\rho^\delta$ be the number of singular values of
$A_\rho(b)$ strictly exceeding $\sqrt2\delta$. Then every mixed-start
controller of uniform error at most $\delta$ uses at least
\begin{equation}\label{eq:r76-robust-cover}
 \min\left\{\sum_H q_H[G:H]:q_H\in\mathbb Z_{\ge0},\quad
       \sum_Hq_H\ge1,\quad
       \sum_Hq_Hm_\rho(H)\ge r_\rho^\delta\ (\rho\in\widehat G)
       \right\}
\end{equation}
labels. The scalar version may be used for binary outputs.

Let $N$ be the number of distinct residual tables $b^h(g)=b(hg)$.
If $N>1$, let
\[
 \Delta_b=\min_{b^h\ne b^k}\max_g
                \operatorname{TV}(b(hg),b(kg))>0.
\]
Whenever $\delta<\Delta_b/2$, a pure-start controller with fewer than
$N$ labels must have a projected mixed preparation. Its total budget
is therefore at least one plus the lower bound
\eqref{eq:r76-robust-cover}.
\end{theorem}
\begin{proof}
Put $a(g)=\lambda P_gD$ and $z_c(g)=b(g,c)-a(g,c)$.
Since $\sum_cz_c(g)=0$ and $\tfrac12\sum_c|z_c(g)|\le\delta$,
the sums of the positive and negative parts are each at most $\delta$,
and $\sum_c|z_c(g)|^2\le2\delta^2$. Unitarity gives
\[
 \left\|\begin{pmatrix}z_{c_1}(g)\rho(g^{-1})\\
                  \vdots\\z_{c_{|C|}}(g)\rho(g^{-1})\end{pmatrix}
       \right\|_{\rm op}
 =\left(\sum_c|z_c(g)|^2\right)^{1/2}\le\sqrt2\delta.
\]
Average over $g$ to obtain
$\|A_\rho(b)-A_\rho(a)\|_{\rm op}\le\sqrt2\delta$.
The Fourier-multiplicity lemma gives
$\operatorname{rank}A_\rho(a)\le m_\rho(X)$.
For completeness, if $B$ has rank at most $m$, its nullspace
intersects the span of the first $m+1$ right singular vectors of $A$
in a nonzero vector. On a unit vector in that intersection,
$\|(A-B)v\|=\|Av\|\ge\sigma_{m+1}(A)$.
This proves \eqref{eq:r76-fourier-robust}; the binary version follows
by averaging the scalar error, whose absolute value is at most
$\delta$, against unitary matrices. Summing orbit multiplicities and
using the mixed-start normal form gives the integer bound.

For the last assertion, suppose the projected preparation were a point
$x$. Its distinct residuals are indexed by points $x\cdot h$ in its
orbit. If two target residuals $b^h,b^k$ were represented by the same
point, the two approximating residuals would coincide for every future
$g$, and the triangle inequality would give distance at most $2\delta$
between $b^h$ and $b^k$. This contradicts the definition of $\Delta_b$.
The orbit must therefore contain at least $N$ points. With fewer than
$N$ total labels this is impossible, so the projected start is mixed
and its extra initializer is charged by the normal-form theorem.
\end{proof}

\begin{corollary}[Uniform error cannot replace a prime orbit]
\label{cor:r76-prime-error}
Let $G=C_p$ with $p$ prime and let $b:G\to\Delta(C)$ be any target.
Put $r(b)=\min_{q\in\Delta(C)}\max_g\operatorname{TV}(b(g),q)$.
For either start convention,
\[
 \mathcal E_M(b)=
 \begin{cases}r(b),&1\le M<p,\\0,&M\ge p.\end{cases}
\]
For binary outputs $r(b)=(\max_g\beta(g)-\min_g\beta(g))/2$.
\end{corollary}
\begin{proof}
Every orbit of $C_p$ has size one or $p$. Thus any action with fewer
than $p$ points is trivial and gives a constant response, regardless
of preparation. A one-label controller attains $r(b)$. On $p$ points
the regular action with a point preparation and decoder $b$ is exact.
For binary outputs the midpoint of the smallest interval containing
all values of $\beta$ is a Chebyshev center of that interval.
\end{proof}

For the symmetric-group target \eqref{eq:r75-symmetric-return}, the
scalar standard block is $tI/[n(n-1)]$ and $\Delta_{b_t}=2t/n$.
Indeed a permutation moving $k$ points changes a residual by at most
$kt/n$, with equality at an inverse representative; the smallest
nonzero $k$ is two. The subgroup-index estimate
\eqref{eq:r75-index-estimate} and the preceding theorem imply
\begin{equation}\label{eq:r76-symmetric-robust-lower}
 \mathcal E_M^{\mathrm p}(b_t)\ge\frac{t}{n(n-1)}
       \qquad(n\ge5,\ 0<t\le1,\ M\le n(n-1)).
\end{equation}
If the error were smaller, the full standard multiplicity $n-1$
would be forced, and the small-error initialization charge would give
at least $n(n-1)+1$ labels; here $n!>n(n-1)$.
This explicit lower bound is not asserted to be the optimal error
constant. At every exact resource jump the finite optimization also
gives a strictly positive robustness radius, without a claimed closed
formula for that radius.

\section{Shared preparation, invariant polytopes, and noncyclic memory}
\label{sec:r77-shared-preparation}

A device that reproduces one prepared response need not reproduce a family
of preparations with the same memory, transitions, and decoder.  We now
solve the structural problem with this common-device requirement.  The
preparation is selected before the command word; it is not revealed to the
terminal decoder.  In particular, one may not choose a different decoder
after learning which preparation was selected.  This convention is what
makes joint realization a resource question.

Fix a finite group $G$, a surjective command homomorphism
$\gamma:A^*\to G$, a finite outcome set $C$, and a nonempty finite
preparation set $I$.  A \emph{shared controller} consists of one family of
stochastic matrices $Q_a$ and one stochastic decoder $D_0$, together with
an initial row for each $i\in I$.  In the pure convention these rows are
prescribed label rows $e_{o_i}$; different preparations may use the same
label when their required laws agree.  All labels, including initialization
labels, are counted.  In the mixed convention the initial rows may be
arbitrary probability rows and their external preparation is not counted.
The result concerns this explicit distinction, not the bit cost of selecting
$i$ or an uncounted label made available to the decoder.

For each $i,g$ let $K_{i,g}\subseteq\Delta(C)$ be closed and convex,
possibly empty.  The requirement is
\begin{equation}\label{eq:r77-shared-specification}
 e_{o_i}Q_wD_0\in K_{i,\gamma(w)}
 \qquad(i\in I,\ w\in A^*),
\end{equation}
with arbitrary initial rows in the mixed convention.  For probability
rows $\Lambda=(\lambda_i)_{i\in I}$ on a finite set put
\[
 \iota(\Lambda)=
 \bigl|\{\lambda_i:i\in I,\ \lambda_i\hbox{ is not a point mass}\}\bigr|.
\]
Repeated rows are counted once.

\begin{theorem}[Shared normal form and exact initialization charge]
\label{thm:r77-shared-normal-form}
The specification \eqref{eq:r77-shared-specification} has a pure shared
controller on at most $M$ labels if and only if there are a finite
$G$-set $X$, probability rows $\lambda_i$ on $X$, and one stochastic
decoder $D$ such that
\begin{equation}\label{eq:r77-shared-core}
 \lambda_iP_gD\in K_{i,g}\quad(i\in I,g\in G),
 \qquad |X|+\iota(\Lambda)\le M.
\end{equation}
For mixed preparations the criterion is $|X|\le M$.  The action may have
a kernel.  Thus a single permutation core suffices for all preparations,
and each distinct mixed preparation costs exactly one additional label
in this normal form.
\end{theorem}
\begin{proof}
Use the identity-word construction in the proof of
Theorem~\ref{thm:r76-convex-normal-form}.  Its averaging matrix $T$,
stochastic idempotent $E=\lim_nT^n$, and disjoint closed classes
$C_1,\ldots,C_k$ depend only on the command matrices, not on the initial
row.  Write $\mu_x$ for their stationary probability rows.  The sandwiched
maps $EQ_uE$ permute these rows and give one action $P_{\gamma(u)}$ on
$X=\{1,\ldots,k\}$.  Put
\[
 e_{o_i}E=\sum_{x\in X}\lambda_i(x)\mu_x,
 \qquad D(x,\cdot)=\mu_xD_0.
\]
The same limit as in that theorem gives
\[
 \lambda_iP_gD
   =\lim_{n\to\infty}e_{o_i}T^nQ_uT^nD_0\in K_{i,g},
 \qquad \gamma(u)=g,
\]
because each approximant is a convex combination of laws for words
representing $g$.  There is no averaging between different preparations.

If $o_i$ lies in a closed class, then $e_{o_i}E$ is its stationary row,
so $\lambda_i$ is a point mass.  Distinct mixed $\lambda_i$ must therefore
arise from distinct initial labels outside the closed classes.  These
classes use at least $k$ labels, so
$\iota(\Lambda)\le M-k$.  This argument also permits an outside initial
label to project to a point mass, and it does not charge that label when
the projected realization no longer needs it.  It permits several initial
labels to have the same projected mixed row, in which case they are merged.
For mixed initial rows the same projection gives $k\le M$ without the
additional count.

Conversely, retain the permutation labels $X$.  For every distinct
nonpoint row $\lambda$ among the $\lambda_i$, adjoin one label $\star_\lambda$
with decoder $\lambda D$ and, on command $a$, transition row
$\lambda P_{\gamma(a)}$ into $X$.  No transition returns to an adjoined
label.  A point row uses its own point of $X$ as initial label; a nonpoint
row uses $\star_\lambda$.  The empty-word law is $\lambda_iD$ and the law
after every nonempty word is $\lambda_iP_{\gamma(w)}D$.  This establishes
\eqref{eq:r77-shared-core} with the exact displayed count.  The mixed
case starts directly from $\lambda_i$ on $X$.
\end{proof}

For exact targets write $\mathcal B=\{b_i:i\in I\}$, where
$b_i:G\to\Delta(C)$; duplicate target tables are identified.
Let $\kappa_G(\mathcal B)$ and $m_G(\mathcal B)$ denote the mixed and pure
shared minima.  The first equals the least $|X|$ in a simultaneous
factorization $b_i(g)=\lambda_iP_gD$; the second includes
$\iota(\Lambda)$.  In particular
\[
 \max_i\kappa_G(b_i)\le\kappa_G(\mathcal B),
 \qquad
 \kappa_G(\mathcal B)\le m_G(\mathcal B).
\]
Neither inequality compares a collection of separately optimized decoders
with a single common decoder.

\subsection{A geometric formula with all initialization labels counted}
Let $\mathcal P_G(\mathcal B)$ be the set of nonempty polytopes in the
finite-dimensional product simplex
\[
 \mathcal Q=\{q:G\to\Delta(C)\}
\]
that contain $\mathcal B$ and are invariant under
$(q\cdot h)(g)=q(hg)$.  The whole product simplex is one such polytope.
Write $v(P)=|\operatorname{vert}P|$.  The polytope is allowed to extend
beyond the linear span of the prescribed responses.  Forbidding this
extension would impose an additional positive-realization constraint.

\begin{theorem}[Invariant-polytope resource formulas]
\label{thm:r77-shared-polytope}
For every nonempty finite target family,
\begin{align}
 \kappa_G(\mathcal B)
   &=\min_{P\in\mathcal P_G(\mathcal B)}v(P),
       \label{eq:r77-polytope-mixed}\\
 m_G(\mathcal B)
   &=\min_{P\in\mathcal P_G(\mathcal B)}
       \left(v(P)+|\mathcal B\setminus\operatorname{vert}P|\right).
       \label{eq:r77-polytope-pure}
\end{align}
The minima are attained.  Every minimizing polytope constructs an optimal
controller: its vertices are permutation labels, and each distinct
prescribed nonvertex response receives one transient initialization label.
\end{theorem}
\begin{proof}
Given a simultaneous factorization, associate to each core label the
profile $q_x(g)=e_xP_gD$.  The set of profiles is invariant under $G$,
and $P=\operatorname{conv}\{q_x:x\in X\}$ belongs to
$\mathcal P_G(\mathcal B)$.  Every vertex is a profile, so $v(P)\le|X|$.
This proves the necessary mixed bound.

For the pure bound, apply the preceding normal form to an arbitrary
$M$-label pure realization.  Let $s=\iota(\Lambda)$ and $k=|X|$, so
$k+s\le M$.  If a target $b_i$ is not a vertex of $P$, it either equals
a nonvertex core profile $q_x$, or is not any core profile.  There are
at most $k-v(P)$ distinct profiles of the first kind: at least $v(P)$
core labels are required to represent the different vertices.  Every
distinct target of the second kind requires a distinct nonpoint
$\lambda_i$, so there are at most $s$ such targets.  Therefore
\[
 v(P)+|\mathcal B\setminus\operatorname{vert}P|
       \le v(P)+(k-v(P))+s\le M.
\]
This is the step that prevents a nonextreme continuing label used by
another preparation from being charged twice.

Conversely, a finite group preserving $P$ permutes its vertices.  Use
these vertices as $X$, with $D(q,c)=q(e,c)$.  Barycentric coordinates
of $b_i\in P$ supply a probability row $\lambda_i$ with
$\lambda_iP_gD=b_i(g)$.  If $b_i$ is a vertex choose its point row.
If it is not a vertex, every barycentric representation is nonpoint.
Choose the same representation for identical target tables.  Different
target tables have different representations, since their barycentric
images differ.  The normal-form construction thus costs exactly
$v(P)+|\mathcal B\setminus\operatorname{vert}P|$ in the pure convention,
and $v(P)$ in the mixed convention.

The feasible integer costs form nonempty subsets of the positive
integers.  Their least elements are realized by some feasible polytope;
the two constructions prove the asserted equalities and attainment.
\end{proof}

Invariant cones in positive realization and the relation between symmetry
and nonnegative factorizations are classical
\cite{Heller65,BenvenutiFarina04,GPT13}.  Formula
\eqref{eq:r77-polytope-mixed} is a response-polytope form of that structure.
The additional resource statement here is the exact shared-device
pure-initialization term in \eqref{eq:r77-polytope-pure}, obtained after
the common identity-word projection.  It neither assumes that a supplied
stochastic controller satisfies the group relations nor identifies linear
rank with positive memory.

\subsection{Joint Fourier multiplicity and exact low-contrast synthesis}
Let
\[
 V_{\mathcal B}=\operatorname{span}_{\mathbb R}
    \{b_i\cdot h:i\in I,h\in G\},\qquad d=\dim V_{\mathcal B}.
\]
For an irreducible unitary representation $\rho$ of degree $d_\rho$,
use the stacked matrix $A_\rho(b_i)$ from
\eqref{eq:r75-fourier-block} and define the \emph{joint} block
\begin{equation}\label{eq:r77-joint-fourier}
 \mathcal A_\rho(\mathcal B)=
   \bigl[A_\rho(b_1)\ \cdots\ A_\rho(b_{|I|})\bigr],
 \qquad r_\rho(\mathcal B)=\operatorname{rank}_{\mathbb C}
                                  \mathcal A_\rho(\mathcal B).
\end{equation}
Unlike the single-preparation rank, $r_\rho(\mathcal B)$ can exceed
$d_\rho$.  The same decoder must accommodate the union of the Fourier
column spaces, rather than their largest separate rank.

\begin{lemma}[Joint multiplicity]
\label{lem:r77-joint-multiplicity}
The multiplicity of $\rho$ in
$V_{\mathcal B}\otimes_{\mathbb R}\mathbb C$ equals
$r_\rho(\mathcal B)$.  Consequently every simultaneous factorization
on a $G$-set $X$ satisfies
\[
 r_\rho(\mathcal B)\le m_\rho(X)\qquad(\rho\in\widehat G).
\]
\end{lemma}
\begin{proof}
Fourier transformation sends the translate $b_i\cdot h$ to
$A_\rho(b_i)\rho(h)$.  The group-algebra decomposition used in
Lemma~\ref{lem:r75-fourier-multiplicity} therefore identifies its
$\rho$-component with
\[
 \left\{\sum_{i\in I}A_\rho(b_i)M_i:
                       M_i\in M_{d_\rho}(\mathbb C)\right\}.
\]
The matrices $M_i$ are independent, since the different generators may
have independent group-algebra coefficients.  This is the space of all
matrices whose columns belong to the sum of the column spaces of the
$A_\rho(b_i)$.  That sum has dimension $r_\rho(\mathcal B)$, so the
right module consists of precisely that many rows of length $d_\rho$.
For a simultaneous positive factorization, the span of all reachable
preparation rows is a submodule of $\mathbb R^X$, and evaluation against
the common decoder maps it onto $V_{\mathcal B}$.  Complete reducibility
therefore embeds $V_{\mathcal B}$ into $\mathbb R^X$, proving the bound.
\end{proof}

Write $\nu_G(V_{\mathcal B})$ for the smallest size of a permutation
module containing $V_{\mathcal B}$.  For one representative $H$ of each
subgroup conjugacy class, Frobenius reciprocity gives
\begin{equation}\label{eq:r77-joint-module-cover}
 \nu_G(V_{\mathcal B})=
 \min\left\{\sum_Hq_H[G:H]:q_H\in\mathbb Z_{\ge0},\quad
       \sum_Hq_Hm_\rho(H)\ge r_\rho(\mathcal B)
                           \ (\rho\in\widehat G)\right\}.
\end{equation}
The equality follows by orbit decomposition and the real-descent part
of Lemma~\ref{lem:r75-fourier-multiplicity}; it applies also to complex
and quaternionic representation types.

\begin{theorem}[Exact noncyclic memory at small common contrast]
\label{thm:r77-shared-low-contrast}
Suppose that all preparations have the same group-averaged law,
\[
 \frac1{|G|}\sum_{g\in G}b_i(g)=\pi\qquad(i\in I),
\]
and put $b_{i,t}(g)=\pi+t\{b_i(g)-\pi\}$ for $0\le t\le1$.
There is an explicitly positive $t_0$ such that
\begin{equation}\label{eq:r77-shared-low-contrast}
 \kappa_G(\{b_{i,t}:i\in I\})
      =\nu_G(V_{\mathcal B})\qquad(0<t\le t_0).
\end{equation}
For every $t>0$ the joint module and its Fourier multiplicities are
unchanged.  The exact pure resource is the invariant-polytope formula
\eqref{eq:r77-polytope-pure}; in particular, if $r=|\mathcal B|$ and
$N=|\{b_i\cdot g:i\in I,g\in G\}|$, then on this interval
\[
 \nu_G(V_{\mathcal B})\le m_G(\{b_{i,t}\})
       \le\min\{N,\nu_G(V_{\mathcal B})+r\}.
\]
At $t=0$ both resources equal one.
\end{theorem}
\begin{proof}
Remove outcomes with $\pi_c=0$; positivity makes them identically zero
for every preparation and group element.  They will be restored as
zero decoder columns.  Let $\overline b$ be the constant table $\pi$.
Group averaging maps every generating residual to $\overline b$, hence
\[
 V_{\mathcal B}=\mathbb R\overline b\oplus W,
 \qquad W=\ker\left(|G|^{-1}\sum_{g\in G}L_g\right).
\]
The nontrivial part $W$ has no invariant vector and all its elements have
zero total output mass.  If $W=0$, all targets equal $\overline b$ and
the assertion is immediate.  Otherwise choose a minimum $G$-set $X$
with $k=\nu_G(V_{\mathcal B})$ and a real equivariant injection
$j:W\to\mathbb R^X$.  Let $F:\mathbb R^X\to W$ be orthogonal
projection onto $j(W)$ followed by $j^{-1}$.  Since the permutation
inner product is invariant, $F$ is equivariant and annihilates every
invariant row.  Write
\[
 p=(1/k,\ldots,1/k),\qquad
 a_i=j(b_i-\overline b),\qquad
 T(x,c)=F(e_x)(e,c).
\]
Then $a_i\mathbf1=0$, $T\mathbf1=0$, $pT=0$, and
$a_iP_gT=b_i(g)-\pi$.  Put
\[
 A=\max_i\|a_i\|_\infty>0,\quad
 \epsilon=\frac1{2kA},\quad
 B=\max_{x,c}\frac{|T(x,c)|}{\pi_c}>0,\quad
 t_0=\min\{1,\epsilon/B\}.
\]
The rows $\lambda_i=p+\epsilon a_i$ have all entries at least
$1/(2k)$.  For $0\le t\le t_0$ the \emph{common} matrix
$D_t(x,c)=\pi_c+(t/\epsilon)T(x,c)$ is stochastic, and
\[
 \lambda_iP_gD_t=\pi+t\{b_i(g)-\pi\}=b_{i,t}(g)
\]
for every preparation simultaneously.  This proves the matching upper
bound, with one decoder and one action.

The equivariant map that is the identity on $\mathbb R\overline b$
and multiplication by $t$ on $W$ is invertible for $t>0$ and sends all
residuals to their contrasted versions.  Thus their joint module and
all equality relations are unchanged.  Lemma~\ref{lem:r77-joint-multiplicity}
and \eqref{eq:r77-joint-module-cover} give the lower bound.  The constructed
preparations have at most $r$ distinct nonpoint rows, so their pure
initialization costs at most $k+r$.  Alternatively, the union of all
residuals is a deterministic shared machine with $N$ labels.  At zero
contrast one common constant decoder and one fixed label suffice.
\end{proof}

The common-mean hypothesis has mathematical content.  Distinct constant
preparation laws can already have positive-factorization cost larger
than their linear rank; contrast does not change them.  For example,
with the trivial group, the four preparation rows
\[
 \frac12\begin{pmatrix}
 1&1&0&0\\0&1&1&0\\0&0&1&1\\1&0&0&1
 \end{pmatrix}
\]
have linear rank three but require four shared mixed labels.  In any
nonnegative rank-one summand, at most one of the four positive diagonal
entries can be covered: covering two would make both corresponding
cross entries positive, whereas at least one is zero.  Four summands
are therefore necessary, and the displayed rows supply a four-label
realization.  This also explains why a separate scalar module bound
for each preparation cannot replace joint positive synthesis.

\subsection{Robust joint certificates and sharp resource separations}
For weights $\alpha_i\ge0$ with $\sum_i\alpha_i=1$, let
\[
 \mathcal A_{\rho,\alpha}(\mathcal B)
   =[\sqrt{\alpha_1}A_\rho(b_1)\ \cdots\
                       \sqrt{\alpha_{|I|}}A_\rho(b_{|I|})].
\]

\begin{theorem}[Joint robust multiplicity certificate]
\label{thm:r77-joint-robust}
If one mixed shared controller approximates every $b_i(\gamma(w))$
to total-variation error at most $\delta$ for every word $w$, there
is a $G$-set $X$ of size at most its memory such that, for every
irreducible $\rho$ and every such weight vector,
\begin{equation}\label{eq:r77-weighted-fourier}
 \sigma_{m_\rho(X)+1}
       (\mathcal A_{\rho,\alpha}(\mathcal B))\le\sqrt2\delta.
\end{equation}
Hence the integer cover \eqref{eq:r77-joint-module-cover}, with each
$r_\rho$ replaced by the number of singular values strictly above
$\sqrt2\delta$ in any chosen weighted block, is a memory lower bound.
Weights may be chosen separately for different irreducibles.
The same lower bound applies to pure shared controllers.
\end{theorem}
\begin{proof}
Apply Theorem~\ref{thm:r77-shared-normal-form} to the closed convex
total-variation balls.  It supplies a common permutation action and
tables $a_i(g)=\lambda_iP_gD$ with the same error bound.  The calculation
in Theorem~\ref{thm:r76-robust-multiplicity} gives
$\|A_\rho(b_i)-A_\rho(a_i)\|_{\rm op}\le\sqrt2\delta$.
For block rows $z=(z_i)_i$, Cauchy--Schwarz therefore gives
\[
 \left\|\sum_i\sqrt{\alpha_i}
       \{A_\rho(b_i)-A_\rho(a_i)\}z_i\right\|
 \le\sqrt2\delta\sum_i\sqrt{\alpha_i}\|z_i\|
 \le\sqrt2\delta\left(\sum_i\|z_i\|^2\right)^{1/2}.
\]
The weighted joint error has operator norm at most $\sqrt2\delta$.
The approximating joint block has rank at most $m_\rho(X)$ by
Lemma~\ref{lem:r77-joint-multiplicity}; zero weights can only reduce
that rank.  The singular-value nullspace argument used in
Theorem~\ref{thm:r76-robust-multiplicity} proves
\eqref{eq:r77-weighted-fourier}.  Orbit decomposition gives the stated
cover.  The pure normal form has a core no larger than its total memory.
\end{proof}

\begin{proposition}[Unbounded joint multiplicity for a fixed group]
\label{prop:r77-shared-involution}
For every integer $q\ge1$ there is a $q$-preparation experiment on
$G=C_2$ with $q+1$ terminal outcomes such that each preparation alone
has pure and mixed memory two, while their common-device minima are
\[
 \kappa_{C_2}(\mathcal B)=m_{C_2}(\mathcal B)=2q.
\]
Its joint sign multiplicity is $q$, although the sign representation
has degree one.  All target response probabilities may be chosen rational and strictly
positive.  Moreover, in the explicit family below every common
approximation with uniform total-variation error
$\delta<\epsilon/\sqrt{2q}$ still requires at least $2q$ labels.
\end{proposition}
\begin{proof}
Take $C=\{0,1,\ldots,q\}$, the uniform law
$\pi=(q+1)^{-1}\mathbf1$, and $v_i=e_i-e_0$ for $1\le i\le q$.
For any rational $0<\epsilon<1/(q+1)$ put
\[
 b_i(e)=\pi+\epsilon v_i,\qquad
 b_i(s)=\pi-\epsilon v_i,
\]
where $s$ is the nonidentity element of $C_2$.  Each individual response
is nonconstant and is realized by a two-cycle with these two decoder
rows, so its exact memory is two.  For a shared realization the joint
sign block is the $(q+1)\times q$ matrix
$\epsilon[v_1\ \cdots\ v_q]$, of rank $q$.  Every two-cycle in a
$C_2$-set contributes one sign copy; fixed points contribute none.
At least $q$ two-cycles, or $2q$ labels, are necessary.  The disjoint
union of the $q$ displayed two-cycles attains this bound, with pure
preparation $i$ at the positive point of its own cycle.

For equal weights $\alpha_i=1/q$, the squared nonzero singular values
of the joint sign block are the eigenvalues of
$(\epsilon^2/q)(I+\mathbf1\mathbf1^{\mathsf T})$.
Its least singular value is at least $\epsilon/\sqrt q$ (also for
$q=1$).  When $\sqrt2\delta<\epsilon/\sqrt q$, the joint robust
certificate still forces $q$ sign copies.  This lower bound allows
arbitrary stochastic transitions and arbitrary stochastic terminal
decoders satisfying the all-word approximation requirement.
\end{proof}

\begin{proposition}[An unbounded pure initialization cost]
\label{prop:r77-many-initializers}
For every integer $r\ge1$ there is an $r$-preparation binary experiment
on $C_2$ with
\[
 \kappa_{C_2}(\mathcal B)=2,
 \qquad m_{C_2}(\mathcal B)=r+1,
\]
although every preparation separately has pure memory two.
\end{proposition}
\begin{proof}
Choose $0<t_1<\cdots<t_r<1/2$ and set
\[
 b_i(e)=(1/2+t_i,1/2-t_i),\qquad
 b_i(s)=(1/2-t_i,1/2+t_i).
\]
Use one two-cycle with decoder first coordinates $1/2+t_r$ and
$1/2-t_r$.  Appropriate mixtures of its two states realize every
preparation, proving the mixed upper two; nonconstancy proves its lower.
The last preparation uses the positive core point, and the other $r-1$
preparations use distinct transient initialization labels.  This gives
$r+1$ pure labels by Theorem~\ref{thm:r77-shared-normal-form}.

For the lower bound, the $r$ distinct empty-word laws require $r$
distinct pure initial labels.  If those were all the labels, every
decoder row would have first coordinate strictly greater than $1/2$.
Every stochastic transition from any of them would then still give a
first-coordinate probability strictly greater than $1/2$.  This
contradicts the required response after $s$, whose first coordinate is
strictly less than $1/2$.  At least one further label is necessary.
\end{proof}

Both families are ordinary finite causal experiments, and the constructions
also admit diagonal quantum realizations: place the decoder laws on the
diagonal, use permutation unitaries for commands, and use the stated
diagonal preparations.  No quantum advantage is asserted.  Their purpose
is to separate simultaneous preparation, noncyclic multiplicity, and
counted pure initialization from the already solved single-preparation
resource.  For arbitrary different mean laws, arbitrary quantum preparation
sets, or nonreversible commands, the displayed common-mean equality is
not asserted.

\section{Finite memory certificates without physical decoder assumptions}
\label{sec:r77-unrestricted-approximation}

The physical-dictionary lower bound uses the geometry of density matrices
at individual hidden labels.  An arbitrary classical simulator need not
have that geometry.  We now give a different certificate, valid for
arbitrary classical decoder tables.  Its conclusion is quantitatively
weaker, but applies to the original algebraic qubit instrument without
an unproved Diophantine hypothesis.  More generally it applies whenever
experimentally accessible preparations have separated exposing tests.

Nonnegative factorizations, support antichains and their relation to
memory and extension complexity are classical.  In particular the exact
zero-diagonal support bound is already in the literature on nonnegative
rank; see \cite[Section 4.1 and Corollary 4.13]{GPT13} and
\cite[Section 3.1]{VandaeleGillisGlineur17}.  The argument below supplies
an explicit error margin and derives its hypotheses from joint stopped
records.  It does not treat an approximate zero as an exact zero, or
assume that a rare conditional distribution inherits the unconditional
approximation error.

\subsection{A robust support certificate}

For $J\ge2$ define the antichain threshold
\begin{equation}\label{eq:r77-antichain-threshold}
 a(J)=\min\left\{m\ge1:
                   \binom{m}{\lfloor m/2\rfloor}\ge J\right\}.
\end{equation}
In particular $\lceil\log_2J\rceil\le a(J)\le J$.

\begin{lemma}[Robust antichain factorization certificate]
\label{lem:r77-robust-antichain}
Let $\lambda_1,\ldots,\lambda_J$ be probability rows on $M$ labels,
and let $d_1,\ldots,d_J\in[0,1]^M$ be arbitrary columns.  If
\begin{equation}\label{eq:r77-antichain-hypotheses}
 \lambda_n d_n\le\epsilon,
 \qquad \lambda_n d_m\ge b\quad(n\ne m),
 \qquad 0<b\le1,
 \qquad \epsilon<\frac{b^2}{2M^2},
\end{equation}
then $J\le\binom{M}{\lfloor M/2\rfloor}$.
The conclusion uses no linearity, noncontextuality or quantum
representation of the decoder columns.
\end{lemma}
\begin{proof}
Put $S_n=\{i:\lambda_{n,i}\ge b/(2M)\}$.  For $n\ne m$, the
nonnegative sum $\lambda_nd_m\ge b$ has a summand
$\lambda_{n,i}d_m(i)\ge b/M$.  Since both factors are at most one,
\[
 \lambda_{n,i}\ge b/M,\qquad d_m(i)\ge b/M.
\]
Thus $i\in S_n$.  On the other hand
\[
 \lambda_{m,i}\le\epsilon/d_m(i)
                \le M\epsilon/b<b/(2M),
\]
so $i\notin S_m$.  Every ordered pair of distinct sets has this
witness, and the $S_n$ form an antichain in the subsets of $[M]$.

For completeness, choose a uniformly random permutation of $[M]$.
For a fixed set $S$ of size $k$, the probability that its elements
are the first $k$ elements of this permutation, in any order, is
$1/\binom{M}{k}$.  The corresponding events for members of an
antichain are disjoint, since the initial sets of one permutation
form a chain.  Hence
\[
 \frac{J}{\binom{M}{\lfloor M/2\rfloor}}
 \le\sum_{n=1}^J\binom{M}{|S_n|}^{-1}\le1,
\]
as required.
\end{proof}

The combinatorial conclusion is sharp at the level of these hypotheses.
For $M\ge2$, let $k=\lfloor M/2\rfloor$, index preparations and
tests by all $k$-element subsets $S$ of $[M]$, and put
\[
 \lambda_S(i)=\frac{\mathbf1_{\{i\in S\}}}{k},
 \qquad d_T(i)=\mathbf1_{\{i\notin T\}}.
\]
Then $\lambda_Sd_S=0$ and $\lambda_Sd_T\ge1/k$ for $S\ne T$, with
exactly $\binom{M}{k}$ preparations.  This construction concerns the
abstract preparation--test table.  It does not construct a small
controller for the qubit's chronological command--report dynamics.

\subsection{An operational certificate from joint records}

Consider any causal experiment with $J$ preparation procedures.  The
$n$th procedure either declares failure or ends at an observable herald
$H_n$.  Only after that herald is a binary terminal test $m\in[J]$
requested.  The choice of $m$, including a policy's future test code,
is not an input to the simulator before the herald.  Its terminal decoder
is one fixed table for each requested
test, independent of the earlier preparation procedure except through
the simulator's current counted label.  These are the usual causal
finite-memory conventions; allowing a decoder to inspect an uncounted
copy of the entire past would define a different resource.

Let $p_n$ be the target probability of $H_n$, and let $f_{nm}$ be the
target rejection probability conditional on $H_n$ and test $m$.
Suppose that every one of these preparation--test experiments has
target expected cost at most $B$, and that
\begin{equation}\label{eq:r77-operational-exposure}
 p_n\ge p_*>0,\qquad f_{nn}=0,\qquad
 f_{nm}\ge s>0\quad(n\ne m).
\end{equation}
The terminal tests and preparation procedures may be classical or
quantum; only their observable probabilities enter the next theorem.

\begin{theorem}[Memory lower bound from approximately exposing tests]
\label{thm:r77-operational-antichain}
Suppose \eqref{eq:r77-operational-exposure} holds, with $J\ge2$ and
$0<s\le1$.  An $M$-label classical simulator that approximates the
complete stopped record of every specified experiment to total-variation
error at most $\delta$ satisfies
\begin{equation}\label{eq:r77-budgeted-antichain}
 \delta\le\frac{p_*s^2}{32M^2}
 \quad\Longrightarrow\quad
 J\le\binom{M}{\lfloor M/2\rfloor}.
\end{equation}
Consequently
\begin{equation}\label{eq:r77-operational-lower}
 \delta\le\frac{p_*s^2}{32J^2}
 \quad\Longrightarrow\quad M\ge a(J).
\end{equation}
Both assertions allow freely mixed or counted pure starts and arbitrary
stochastic terminal decoder tables.  They require only the finite
collection of $J^2$ displayed preparation--test experiments.
\end{theorem}
\begin{proof}
For the $n$th procedure let $q_n$ be the simulator's herald probability,
and let $\lambda_n$ be its conditional distribution on its counted
labels at that herald.  It does not depend on the later requested $m$.
Write $d_m(i)\in[0,1]$ for the rejection probability of test $m$ at
label $i$.  Total variation controls both the herald event and the
joint herald-and-rejection event, so
\begin{equation}\label{eq:r77-joint-error}
 |q_n-p_n|\le\delta,\qquad
 |q_n\lambda_nd_m-p_nf_{nm}|\le\delta.
\end{equation}
Under the hypothesis in \eqref{eq:r77-budgeted-antichain},
$\delta/p_n\le s^2/(32M^2)\le1/32$, so $q_n>0$ and
\[
 \lambda_nd_n\le\frac{\delta}{p_n-\delta}
                 \le\frac{2\delta}{p_*}
                 \le\frac{s^2}{16M^2}.
\]
For $m\ne n$ put $t=\delta/p_n\le s^2/(32M^2)\le s/4$.
Equation~\eqref{eq:r77-joint-error} gives
\[
 \lambda_nd_m\ge\frac{p_ns-\delta}{p_n+\delta}
                =\frac{s-t}{1+t}\ge\frac{s}{2}.
\]
Apply Lemma~\ref{lem:r77-robust-antichain} with $b=s/2$ and
$\epsilon=s^2/(16M^2)<b^2/(2M^2)$.  This proves
\eqref{eq:r77-budgeted-antichain}.  If $M\ge J$, then $M\ge a(J)$
already.  Otherwise the hypothesis in \eqref{eq:r77-operational-lower}
implies the one in \eqref{eq:r77-budgeted-antichain}, and gives the
same conclusion.  No conditional error estimate was presumed before
\eqref{eq:r77-joint-error}; normalization was performed only after
the herald mass was controlled.
\end{proof}

This theorem keeps the cost budget in the target law.  It makes no
assertion about the simulator's expected duration.  The hidden label
can have any classical interpretation, and the decoder for one binary
POVM can be unrelated to decoders for other POVMs.

\subsection{An unrestricted lower bound for the original algebraic instrument}

For the undamped qubit instrument
\eqref{eq:r75-noisy-instrument}, fix one commanded Bloch angle
$\theta$ and put $z=e^{i\theta}$ and $\omega=\theta/(2\pi)$.  Let
$N_{\rm cl}^{\rm m}(\delta;B)$ and
$N_{\rm cl}^{\rm p}(\delta;B)$ be the least label counts approximating
every common visible stopping and terminal-POVM policy of target expected
calls at most $B$, with arbitrary classical terminal decoder tables.
Transitions are fixed independently of later terminal requests.  Thus
these classes include the physical dictionaries of
Section~\ref{sec:r76-physical-approximation}, but impose no Born
representation on their individual labels.

Fix $B>0$ and write $\zeta=\min\{1,7B/8\}$.  For $J\ge2$ and
a phase which is not a root of unity put
\begin{equation}\label{eq:r77-phase-margin}
 s_J=\min_{1\le q<J}\frac{1-\operatorname{Re}(z^q)}{2}>0.
\end{equation}

\begin{theorem}[Unrestricted finite certificates and arithmetic rates]
\label{thm:r77-unrestricted-qubit}
For either start convention, the following statements hold.

\smallskip\noindent
\textup{(i)} For every non-root phase and every integer $J\ge2$,
\begin{equation}\label{eq:r77-phase-memory}
 0<\delta\le\frac{\zeta\,8^{-J}s_J^2}{32J^2}
 \quad\Longrightarrow\quad N_{\rm cl}(\delta;B)\ge a(J).
\end{equation}
This lower bound uses $J$ terminal projectors and finitely many
preparation prefixes, with no assumption on individual decoder rows.

\smallskip\noindent
\textup{(ii)} At the original algebraic phase
$z=(3+4i)/5$, the explicit sufficient condition is
\begin{equation}\label{eq:r77-rational-phase-memory}
 0<\delta\le\frac{25\zeta}{128\cdot200^J J^2}
 \quad\Longrightarrow\quad N_{\rm cl}(\delta;B)\ge a(J).
\end{equation}
No badly-approximable assertion about $\theta/(2\pi)$ is required.

\smallskip\noindent
\textup{(iii)} If $\theta=2\pi\omega$ and
$\|q\omega\|_{\mathbb R/\mathbb Z}\ge c/q$ for every $q\ge1$,
then another explicit sufficient condition is
\begin{equation}\label{eq:r77-badly-approx-memory}
 0<\delta\le\frac{\zeta c^4}{2\cdot8^J J^6}
 \quad\Longrightarrow\quad N_{\rm cl}(\delta;B)\ge a(J).
\end{equation}

\smallskip\noindent
In case \textup{(ii)}, and also in case \textup{(iii)},
\begin{equation}\label{eq:r77-unrestricted-growth}
 \liminf_{\delta\downarrow0}
 \frac{N_{\rm cl}^{\rm m}(\delta;B)}
      {\log_2\log_8(1/\delta)}\ge1,
 \qquad
 \liminf_{\delta\downarrow0}
 \frac{N_{\rm cl}^{\rm p}(\delta;B)}
      {\log_2\log_8(1/\delta)}\ge1.
\end{equation}
The inherited physical constructions give, at the same fixed phase,
$N_{\rm cl}^{\rm m}\le K_\omega+L-1$ and
$N_{\rm cl}^{\rm p}\le K_\omega+L$, where $L$ is
\eqref{eq:r76-qubit-depth}; at $z=(3+4i)/5$ one may take $K_\omega=7$.
The unrestricted lower and upper therefore leave a gap between
logarithmic-in-logarithm and logarithmic growth.  A matching logarithmic
lower is asserted only for the physical-dictionary class under the
hypotheses of Theorem~\ref{thm:r76-physical-logarithm}.
\end{theorem}
\begin{proof}
For $n=0,\ldots,J-1$, begin with probability $\zeta$ and otherwise
stop with no call.  If begun, follow the report prefix $1\,0^n$,
stopping immediately on a different report.  Success is the herald
$H_n$.  The target probability and conditional terminal state are
\[
 p_n=\zeta8^{-(n+1)},\qquad
 \rho_n=U_\theta^n\tau(U_\theta^*)^n.
\]
In the real Bloch plane $\rho_n$ has vector
$(\cos(n\theta),\sin(n\theta))$.  The target expected number of
calls, before the cost-free terminal test, is
\[
 \zeta\sum_{j=0}^{n}8^{-j}\le\zeta\,8/7\le B.
\]
After success request the binary terminal POVM
$(\rho_m,I-\rho_m)$, and call its second outcome rejection.  Its
target conditional rejection probability is
\[
 f_{nm}=1-\tr(\rho_n\rho_m)
       =\frac{1-\cos((n-m)\theta)}2.
\]
Thus $f_{nn}=0$, $f_{nm}\ge s_J$ for $n\ne m$, and
$p_n\ge p_*:=\zeta8^{-J}$.  These are the hypotheses of
Theorem~\ref{thm:r77-operational-antichain}, which gives (i).
The simulator's terminal projector request is made only after the
prefix.  Causality is precisely what gives one common label
distribution $\lambda_n$ for all its $J$ possible terminal tests.

For (ii), write
\[
 z^q=\frac{A_q+iB_q}{5^q},\qquad A_q,B_q\in\mathbb Z.
\]
The phase is not a root of unity: otherwise $z+z^{-1}=6/5$ would
be both a rational number and an algebraic integer, hence an integer.
Consequently $\operatorname{Re}z^q<1$ for every $q\ge1$, and
$5^q-A_q$ is a positive integer.  Therefore
\[
 \frac{1-\operatorname{Re}z^q}{2}
 =\frac{5^q-A_q}{2\cdot5^q}\ge\frac1{2\cdot5^q},
 \qquad s_J\ge\frac1{2\cdot5^{J-1}}.
\]
Substitution in (i) gives
\[
 \frac{\zeta8^{-J}}{32J^2}\frac1{4\cdot25^{J-1}}
 =\frac{25\zeta}{128\cdot200^J J^2},
\]
which proves the claimed rational certificate.  The phase here is
the Bloch rotation, as in \eqref{eq:r75-noisy-instrument}; its unitary
uses the half-angle.  No angle is squared in this calculation.

For (iii), $\sin(\pi t)\ge2t$ for $0\le t\le1/2$ gives,
for $1\le q<J$,
\[
 \frac{1-\cos(2\pi q\omega)}2
 =\sin^2(\pi q\omega)
 \ge4\|q\omega\|_{\mathbb R/\mathbb Z}^2
 \ge\frac{4c^2}{q^2}\ge\frac{4c^2}{J^2}.
\]
Part (i) then yields \eqref{eq:r77-badly-approx-memory}.

For fixed $B$, the sufficient condition in (ii) holds with
$J=\log_{200}(1/\delta)-O_B(\log\log(1/\delta))$ chosen as an
integer slightly below this expression.  The corresponding choice in
(iii) is $J=\log_8(1/\delta)-O_{B,c}(\log\log(1/\delta))$.
Since $a(J)\ge\log_2J$, both imply
\eqref{eq:r77-unrestricted-growth}.  The asserted uppers are the
unchanged orbit-prefix constructions of
Theorem~\ref{thm:r76-physical-logarithm} and its final algebraic-phase
paragraph, which are also admissible in the unrestricted class.
\end{proof}

In fact the exact integer lower $a(J)$ retains a second-order term:
the central-binomial estimate
$\binom{m}{\lfloor m/2\rfloor}=\Theta(2^m/\sqrt m)$ gives
$a(J)=\log_2J+\tfrac12\log_2\log_2J+O(1)$.
This is a refinement of the support certificate, not an assertion of
the optimal asymptotic simulator size.  The lower uses terminal
projector requests, whereas the earlier exact infinite-memory
obstruction already used report words alone.  These are different
sets of operational probes.  No positive extra initializer charge is
deduced for the unrestricted approximation problem.

\begin{corollary}[Every non-root algebraic phase has the same lower order]
\label{cor:r77-algebraic-unrestricted}
Let $z$ be algebraic with $|z|=1$ and not a root of unity.  Let its
degree be $D$, choose an integer $A\ge1$ such that $Az$ is an
algebraic integer, and enumerate its conjugates as
$z=z_1,z_2,\ldots,z_D$.  Choose any supplied real bound
\[
 H\ge\max\left\{2,
             A^D\prod_{j=2}^D\max(1,|z_j|)\right\}.
\]
At that fixed phase, for either start convention,
\begin{equation}\label{eq:r77-algebraic-margin}
 0<\delta\le
 \frac{\zeta H^4}{32\cdot16^D\,(8H^4)^J J^2}
 \quad\Longrightarrow\quad N_{\rm cl}(\delta;B)\ge a(J).
\end{equation}
In particular \eqref{eq:r77-unrestricted-growth} holds.  No
badly-approximable premise is necessary for this lower order.
\end{corollary}
\begin{proof}
The number $A^q(z^q-1)=(Az)^q-A^q$ is a nonzero algebraic integer.
Its norm is a nonzero integer, so
\[
 1\le A^{qD}|z^q-1|\prod_{j=2}^D|z_j^q-1|
   \le 2^{D-1}H^q|z^q-1|.
\]
Because $|z|=1$,
\[
 \frac{1-\operatorname{Re}z^q}{2}
   =\frac{|1-z^q|^2}{4}\ge4^{-D}H^{-2q}.
\]
It follows that $s_J\ge4^{-D}H^{-2(J-1)}$.
Substitute this in \eqref{eq:r77-phase-memory} to obtain
\eqref{eq:r77-algebraic-margin}.  A choice
$J=\log_{8H^4}(1/\delta)-O_{B,H,D}(\log\log(1/\delta))$
then gives the claimed lower order.  This elementary algebraic norm
bound does not supply the stronger $s_J\asymp J^{-2}$ bound used for
the sharp physical-dictionary logarithmic law.
\end{proof}

\section{The exact cost of initializing a stationary response}
\label{sec:r75-stationary-initialization}

The distinction between a freely chosen initial distribution and a
prescribed initial label is not peculiar to a periodic example.  There is
an exact general law.  Its proof is an elementary elimination in the
space of future response functions, within the classical positive
realization framework \cite{Heller65,MonrasWinter14,MonrasWinter16}.  It does not require
hidden states to lie in the minimal linear realization, or require their
unobserved transition matrix to be irreducible.

Let $A$ be a nonempty finite command alphabet, with a finite report set
$Y_a$ for each $a\in A$.  A finite classical realization consists of
nonnegative matrices $T_{a,y}$ such that
\[
 T_a:=\sum_{y\in Y_a}T_{a,y},\qquad T_a\one=\one,
\]
and an initial probability row $\mu$.  For a command--report word
$w=(a_1,y_1)\cdots(a_n,y_n)$ write
$T_w=T_{a_1,y_1}\cdots T_{a_n,y_n}$ and
$p(w)=\mu T_w\one$.  Optional terminal readouts have decoder
columns $d_{b,z}\ge0$, with $\sum_zd_{b,z}=\one$, and response
$p(w;b,z)=\mu T_wd_{b,z}$.  All following equalities include every specified
terminal readout when these are part of the experiment.

Call the prepared response \emph{predictively stationary under $\nu$}
if $\nu$ is a probability row on $A$ with every $\nu(a)>0$ and
\begin{equation}\label{eq:r75-predictive-stationarity}
 \sum_{a\in A}\nu(a)\sum_{y\in Y_a}p((a,y)w)=p(w)
 \qquad\text{for every future word }w.
\end{equation}
When terminal readouts are specified, additionally require
\begin{equation}\label{eq:r75-terminal-stationarity}
 \sum_{a\in A}\nu(a)\sum_{y\in Y_a}p((a,y)w;b,z)=p(w;b,z)
 \qquad\text{for all }w,b,z,
\end{equation}
including the empty word $w$.
This is an equality of observable responses; it does not assume
$\mu\sum_a\nu(a)T_a=\mu$ in a particular hidden realization.  For one
command and no terminal readout it is ordinary stationarity.  A response
has a one-label realization precisely when its report probabilities
factor as $p(w)=\prod_jp_{a_j}(y_j)$, and every terminal readout has a
fixed distribution independent of the preceding word.  For one command
this is the i.i.d. case.

If a stationary hidden row is desired, it can be chosen without adding
labels. Put $\overline T=\sum_a\nu(a)T_a$. Every $\mu\overline T^j$
has the same future response by the displayed stationarity identities.
A subsequential limit of the Ces\`aro averages
$H^{-1}\sum_{j=0}^{H-1}\mu\overline T^j$ is stationary for
$\overline T$ and retains that response. Thus the freely mixed minimum
below agrees with the convention requiring a stationary hidden row for
the averaged command, and with the usual stationary HMM convention when
there is only one command.

Let $n_{\mathrm{mix}}(p)$ be the minimum number of labels with an arbitrary
initial probability row, and let $n_{\mathrm{pure}}(p)$ be the minimum when
one prescribed initial label is counted.  The minima are $+\infty$ when
there is no finite realization.

\begin{theorem}[Stationary initialization law]
\label{thm:r75-stationary-initializer}
Suppose \eqref{eq:r75-predictive-stationarity} holds, together with
\eqref{eq:r75-terminal-stationarity} when terminal readouts are included.
If the response has
a one-label realization, both minima are one.  Otherwise
\begin{equation}\label{eq:r75-initialization-law}
 n_{\mathrm{pure}}(p)=n_{\mathrm{mix}}(p)+1,
 \qquad (+\infty)+1:=+\infty.
\end{equation}
The lower holds for all real nonnegative transition tables, including
reducible and nonminimal linear realizations.  The upper processes the
first actual command while entering the core; it neither discards a
physical call nor inserts an uncounted preliminary randomization.
\end{theorem}
\begin{proof}
Take any finite pure-start realization, with initial label $o$, and let
$f_i$ denote the complete future response from hidden label $i$.  Thus
$f_i(w)=e_i^{\mathsf T}T_w\one$, with the analogous values for terminal
readouts.  Put $\overline T=\sum_a\nu(a)T_a$.  Stationarity gives the
functional identity
\begin{equation}\label{eq:r75-initial-functional-identity}
 f_o=\sum_j\overline T_{oj}f_j.
\end{equation}
If $\overline T_{oo}=1$, nonnegativity and full support of $\nu$ imply
$(T_a)_{oo}=1$ for every $a$.  All report transitions from $o$ then return
to $o$.  Their probabilities, and the terminal readouts at $o$, constitute
a one-label realization.  In the case under consideration this is
excluded.  Hence
\[
 \beta_j=\frac{\overline T_{oj}}{1-\overline T_{oo}}\quad(j\ne o)
\]
is a probability row and $f_o=\sum_{j\ne o}\beta_jf_j$.

Delete label $o$.  On the remaining labels define
\begin{equation}\label{eq:r75-elimination-tables}
 \widehat T_{a,y}(i,j)
   =T_{a,y}(i,j)+T_{a,y}(i,o)\beta_j,
 \qquad i,j\ne o,
\end{equation}
and retain their terminal decoders.  These matrices are nonnegative and
their sums over $y$ are stochastic.  The recursion
$f_i((a,y)w)=\sum_jT_{a,y}(i,j)f_j(w)$, with the functional identity for
$f_o$ substituted on its right side, is exactly the recursion for
\eqref{eq:r75-elimination-tables}.  Induction on word length, starting with
the empty word and each terminal decoder, proves that every retained
$f_i$ is unchanged.  Initial row $\beta$ therefore realizes $f_o$ with
one fewer label.  This proves $n_{\mathrm{mix}}\le
n_{\mathrm{pure}}-1$ whenever a finite pure realization exists.

Conversely, given a $k$-label mixed-start realization $(\mu,T_{a,y})$,
append one label $o$.  Its report transition on command $a$ enters the
core with row $\mu T_{a,y}$ and never returns to $o$; the core retains its
original transitions.  A terminal readout at $o$ has row $\mu d_{b,z}$.
The first physical command is already included in this row, so induction
gives the original law at every word length, including zero.  This uses
$k+1$ labels and proves the reverse inequality.  A finite pure realization
is itself a finite mixed one, so the two finiteness assertions are
equivalent.  The one-label exception is immediate.
\end{proof}

\begin{corollary}[Stationary processes]
\label{cor:r75-stationary-process}
For every stationary finite-alphabet process that is not i.i.d., the
minimum pure-start edge-emitting hidden Markov realization has exactly
one more label than the minimum freely mixed realization, with the same
infinite-value convention.  In particular, the $q$ versus $q+1$ conclusion
of Theorem~\ref{thm:r74-depolarizing-obstruction} is an instance of a
general stationary principle, while its phase-dependent $q$ remains a
separate spectral conclusion.
\end{corollary}

Equality of these complete response functions also preserves the law
under every common randomized policy for commands, stopping and terminal
readout.  For each finite stopped history, multiply the equal response by
the same policy probabilities.  Summation gives all finite stopping
events, and decreasing continuation events give the nontermination
probability.  Consequently all nonnegative costs measurable from that
record have equal expectations, including infinite expectations.  These
are process-realization label counts.  If the external policy has its own
persistent memory, that memory is counted separately in a closed-loop
implementation; it is not supplied free by the theorem.

\section{The causal task-order transfer theorem}\label{sec:central}
Suppose that every terminating state-stopped competitor reaches at least $B_0\ge1$ live labels under $\theta_*$. Such a certificate must concern actual physical command compatibility, not printed action names. Define
\begin{equation}\label{eq:widths}
\begin{array}{c|cc}
 i&m_-(M)&m_+(M)\\\hline
 \cp&M&\lfloor M/D_0\rfloor\\
 \ev&M&\lfloor M/D_0\rfloor\\
 \st&M-B_0&\lfloor(M-B)/D_0\rfloor .
\end{array}
\end{equation}
The event upper also requires $M\ge B$. The trivial state-stopped lower $B_0=1$ is valid for a procedure initially waiting to make a call.

\begin{theorem}[Affine task-order transfer with a paid reference and common construction]\label{thm:transfer}
Assume all the prepared-kernel, full-report, task-order, effort, occupation, resource and common-construction hypotheses of Section~\ref{sec:data}. Fix $N,M$, put $b_-=N/\kappa_-$, $b_+=N/\kappa_+$, and suppose $b_-\ge1$ and $m_-(M)\ge1$. Set
\[
 t_N=(1-A_N)_+/u_-.
\]
If $t_N>0$, every one of the three risks satisfies
\begin{equation}\label{eq:central-lower}
 R_i^\Theta(N,M)\ge
 \bigl[l_-t_N\Q_{m_-(M)}(b_-/t_N)-E_N\bigr]_+.
\end{equation}
If $t_N=0$, the asserted lower is zero, with no evaluation of $\Q$ at infinity. An empty legal class still has risk $+\infty$.

Suppose also $m_+(M)\ge1$, $b_+\ge1$, and, for event stopping, $M\ge B$. If
\begin{equation}\label{eq:beta}
 \beta_N=\frac{l_+b_+}{1+b_+a_+}\ge1,
\end{equation}
one common Borel program with actual expected cost at most $N$ in every model attains
\begin{equation}\label{eq:central-upper}
 R_i^\Theta(N,M)\le
 \frac{u_+(1+b_+a_+)}{l_+}\Q_{m_+(M)}(\beta_N)+b_+e_+.
\end{equation}
It commits almost surely. If $\beta_N<1$, the always-accept first-attempt construction instead gives $R_i^\Theta(N,M)\le u_+\Q_{m_+(M)}(1)+e_+$ at cost at most $\kappa_+\le N$.

The lower is against every legal competing strategy and uses one actual witness. The upper is one parameter-independent strategy for all models. The two effort constants need not agree. Neither direction assumes that the comparison experiment belongs to $\Theta$ or that a normalized predictive update is uniformly contracting.
\end{theorem}

\begin{corollary}[Scale-uniform relative transfer]\label{cor:relative}
Under Theorem~\ref{thm:transfer}, suppose $A_N=E_N=a_+=e_+=0$. If $l_+b_+\ge1$, then
\begin{equation}\label{eq:relative}
 \frac{l_-}{u_-}\Q_{m_-(M)}(u_-b_-)
 \le R_i^\Theta(N,M)
 \le\frac{u_+}{l_+}\Q_{m_+(M)}(l_+b_+).
\end{equation}
The factors $l_\pm,u_\pm$ may be fixed nonunit constants, independent of budget and resolution. In particular no term proportional to $b_\pm$ times a fixed absolute calibration error is required in this relative-order class.
\end{corollary}
\begin{proof}
Substitute $t_N=1/u_-$ and $\beta_N=l_+b_+$ in the theorem.
\end{proof}

\begin{corollary}[Separate exact checkpoint and autonomous laws]\label{cor:exact}
Suppose the certificates are exact: $l_\pm=u_\pm=1$, all defects vanish, $\kappa_-=\kappa_+=\kappa$, and exports are the identity. Then, for $N\ge\kappa$,
\[
 R_\cp^\Theta(N,m)=\Q_m(N/\kappa).
\]
The constructed autonomous upper uses $B+m$ state-stopped labels, or $\max(B,m)$ event-stopped labels. If the mandatory live count is sharp, $B_0=B$, then
\begin{align}
 R_\st^\Theta(N,B+m)&=\Q_m(N/\kappa)\quad(m\ge1),\label{eq:typed-exact}\\
 R_\ev^\Theta(N,M)&=\Q_M(N/\kappa)\quad(M\ge B).\label{eq:event-exact}
\end{align}
The subtraction by $B$ belongs only to the state-stopped model. Nonidentity exports require the expanded dictionary in \eqref{eq:widths}, unless an optimal dictionary is invariant under all exports. These exact stopping distinctions are inherited from the complete companion and are reproved here to fix the resource semantics.
\end{corollary}
\begin{proof}
The same lower and upper profile coincide in each asserted interface. Identity exports give $D_0=1$. The packing and converse widths are proved in Lemma~\ref{lem:roles}. No lower on live states is used to forbid event-stopped reuse.
\end{proof}

The familiar absolute-defect theorem is also a specialization: take $l_\pm=u_\pm=1$ while retaining $A_N,E_N,a_+,e_+$. Thus relative and absolute calibration are composed in one statement. The chosen metric is the common scored task throughout; it is not changed between the converse and the construction.

\section{Proof of the central theorem}
\begin{lemma}[Profile attainment]\label{lem:attainment}
For every finite $m\ge1$ and $b\ge1$, the minimum in \eqref{eq:profile} is attained. The profile is nonincreasing in $m,b$ and left-continuous in $b>1$.
\end{lemma}
\begin{proof}
Write a dictionary as an ordered tuple in the compact set $H^m$, permitting repetitions. The feasible densities are a weak-star closed subset of the radius-$b$ ball of $L^\infty(\nu)$: positivity, the upper bound and the mass equality are each closed under pairing with $L^1(\nu)$. They are therefore weak-star compact. As the tuple varies, its squared-distance function varies uniformly on $H$, since all points belong to a bounded set. Along a convergent product net, split the objective difference into this uniform variation and the weak-star pairing against the fixed bounded measurable loss at the limit tuple. This proves continuity of the objective and attainment. Nets avoid an unwarranted sequential compactness assumption. Monotonicity follows from inclusion of feasible sets. For left continuity take an optimizer $h$ at $b>1$ and let $h_\lambda=(1-\lambda)h+\lambda$. It has unit mass and cap $b_\lambda=(1-\lambda)b+\lambda<b$. Its objective tends to the optimum as $\lambda\downarrow0$, while monotonicity bounds the optimum at $b_\lambda$ between the optimum at $b$ and this objective.
\end{proof}

\begin{lemma}[Actual-history comparison measure]\label{lem:comparison}
Fix any robustly admissible competitor of cost at most $N$, with its fixed terminal dictionary $C$. There are finite nonnegative measures $\xi_c$ on $X$, $c\in C$, such that $\xi=\sum_c\xi_c$ obeys
\begin{align}
 \xi&\le (\E_*A_\tau)\nu\le b_-\nu,\label{eq:measure-cap}\\
 1&\le u_-\xi(X)+A_N,\label{eq:measure-mass}\\
 \E_*\|Z-\widehat r\|^2
 &\ge l_-\sum_{c\in C}\int d_c\,d\xi_c-E_N.
 \label{eq:measure-risk}
\end{align}
The measure $\xi$ is not asserted to be the terminal law of any comparison execution.
\end{lemma}
\begin{proof}
Adjoin an independent uniform random variable to the current report on both sides. This is current-call randomness, not a seed retained between calls; it is sufficient to represent the competitor's randomized output and stopping decision. A finite-output kernel has such a representation by cumulative probabilities. At a fixed reached productive history $h$ and chosen command $a$, the post-report decision then gives an acceptance indicator $q_h(y)$ and an output $r_h(y)\in C$ on acceptance. Give $r_h$ an arbitrary fixed value on rejection. The notation $y$ now includes the current independent random variable. Integrating the original certificates over that variable preserves their constants. Dependence on the entire earlier history is permitted in these functions.

For each Borel $D\subset X$, set
\begin{equation}\label{eq:xi-definition}
 \xi_c(D)=\E_*\sum_{j\le\tau:\ j\ {
m productive}}
 \int\one_{\{\pi_{H_{j-1},a_j}(y)\in D\}}
 q_{H_{j-1}}(y)\one_{\{r_{H_{j-1}}(y)=c\}}
 P^0_{H_{j-1},a_j}(dy,H).
\end{equation}
The outer expectation is the actual pre-call occupation at $\theta_*$. All terms are jointly measurable and nonnegative. Tonelli defines the measures, whose total is finite because $\E_*A_\tau\le b_-$. Summing over $c$ and dropping the acceptance indicator gives \eqref{eq:measure-cap}: each comparison row has mark law $\nu$.

The analogous sum with the actual rows instead of $P^0$ has mass exactly one, because the actual procedure commits exactly once almost surely. Apply \eqref{eq:lower-mass} at each actual reached history, then sum. The occupation certificate yields \eqref{eq:measure-mass}. Apply \eqref{eq:lower-task} in the same way to the disjoint actual commitment events. Their expected losses sum to the actual risk. The comparison contribution is $d_{r_h(y)}(\pi_{h,a}y)$ by \eqref{eq:full-moments}; this identity conditions on the full current report and therefore does not overlook extra information seen by the competitor. The accumulated loss defect is at most $E_N$, giving \eqref{eq:measure-risk}. One may first truncate the number of calls. Monotone convergence for the nonnegative contributions and the finite defect bounds justify the limit.

No actual past is replaced by a comparison past. In particular the competitor need not terminate, or obey a cost bound, if hypothetically run under the comparison experiment. This distinction is why the argument applies when that experiment is not an operating model.
\end{proof}

\begin{remark}[Why the full report is required]
A constant mark can conceal a visible fair bit which also equals the future audit. A comparison using only that mark has $f=0,w=1$, but a two-value decoder reading the bit predicts it with zero loss. The hypotheses therefore condition on the full report seen by the competitor; matching a compressed mark marginal alone cannot establish the lower.
\end{remark}

\begin{lemma}[Live roles, output roles and physical compatibility]\label{lem:roles}
A program with $\ell$ live roles and $t$ output roles has a state-stopped implementation with $\ell+t$ labels and an event-stopped implementation with $\max(\ell,t)$ labels, preserving every path cost and complete stopped scored law. Conversely an event-stopped $M$-label program has at most $M$ decoder values. A state-stopped $M$-label program with at least $B_0$ reached live labels has at most $M-B_0$ terminal decoder values.

For a physical live-state certificate, consider finitely many mandatory pre-call contexts $I$ with no free context input. Suppose the commands legal at a history of context $i$ are exactly $\mathcal U_i$, and every terminating competitor visits every context with positive probability under the same witness. This positive probability may depend on the competitor; no uniform lower bound over competitors is required. Let $b_{\rm phys}$ be the least size of a cover of $I$ by subsets $J$ with $\bigcap_{i\in J}\mathcal U_i\ne\varnothing$. Then at least $b_{\rm phys}$ live labels are required. This certificate is invariant under duplicate names and bijective stateless recodings of the same physical commands.
\end{lemma}
\begin{proof}
For state stopping take disjoint copies of the live and terminal role sets and type them zero and one. For event stopping embed the live and terminal role sets separately into one alphabet of size $\max(\ell,t)$. On continuing transitions use the live embedding; on commitment emit the final event and use the terminal embedding. A command row is read only on a continuing branch, while the decoder is read only after commitment. Sharing their label therefore causes no conflict and uses no continuing mode bit. Couple the original and recoded coins. Every finite stopped prefix and every path cost agree, and almost-sure termination gives equality of the stopped scored laws. A fixed decoder has one value per retained label. In state stopping the global halt map excludes each reached live label from the output set, proving its subtraction. Event stopping has no such exclusion.

A live label reached in contexts $J$ must have one command distribution supported on all the corresponding $\mathcal U_i$. Since $J$ is finite, the union of the exceptional null command sets is null; hence their intersection supports a probability and is nonempty. The sets of contexts served by the labels cover $I$, proving the bound. Conversely a nonempty intersection contains a command and hence a Dirac command distribution. Replacing names without changing executed kernels, legal domains, costs and reports does not change these intersections. A translator which itself remembers a hidden phase is not such a recoding; its memory must be charged.
\end{proof}

\begin{proof}[Proof of Theorem~\ref{thm:transfer}]
Fix a competing program. Lemma~\ref{lem:roles} bounds its decoder width by $m_-(M)$. Let $s=\xi(X)$ in Lemma~\ref{lem:comparison}. If $t_N>0$, then \eqref{eq:measure-mass} gives $s\ge t_N$. Scale the measure to $\xi'=(t_N/s)\xi$. It has mass $t_N$ and remains bounded above by $b_-\nu$. Thus $h=(d\xi'/d\nu)/t_N$ has unit mass and cap $b_-/t_N$. Nearest-dictionary projection and nonnegativity imply
\[
 \sum_{c\in C}\int d_c\,d\xi_c
 \ge\int [w+\dist(f,C)^2]\,d\xi
 \ge t_N\Q_{m_-(M)}(b_-/t_N).
\]
Together with \eqref{eq:measure-risk} this proves the asserted lower for this competitor. Its worst-model risk dominates its risk at $\theta_*$. Taking the infimum over competitors gives \eqref{eq:central-lower}, without exchanging an infimum and a supremum. Nonnegative loss gives the zero truncation and the stated case $t_N=0$.

For the upper assume \eqref{eq:beta}. Choose an attained optimizer $(C,h)$ at width $m_+(M)$ and cap $\beta_N$. Let $r_C(x)$ be the first nearest point to $f(x)$ in a fixed ordering of $C$. After each attempt of $\mathcal G$, accept $x$ with probability $h(x)/\beta_N$, export $T_Jr_C(x)$ and commit; otherwise return to the start routing label. All these operations are Borel. The choice $J$ is used within the current attempt; any part retained thereafter is charged through the live role or exported decoder, not a free seed. The exports are assumed jointly measurable and take at most $D_0|C|$ distinct values, so their finite decoder and transitions are Borel as well.

Conditional on every surviving attempt history, the acceptance probability is at least
\begin{equation}\label{eq:hazard}
 p=\frac{l_+}{\beta_N}-a_+=\frac1{b_+}>0.
\end{equation}
Consequently the probability of reaching attempt $j+1$ is at most $(1-p)^j$, by induction on conditional survival probabilities. In particular commitment is almost sure and the expected number of attempts is at most $1/p=b_+$. Each begun attempt has conditional expected cost at most $\kappa_+$, not merely a cost conditional on a favorable outcome. Summing the predictable attempt costs gives actual expected cost at most $\kappa_+b_+=N$.

The conditional unnormalized loss contributed by an attempt is at most $u_+\Q_{m_+(M)}(\beta_N)/\beta_N+e_+$. Sum its contributions on the disjoint terminal events, using the same survival bound. The result is at most
\[
 b_+\left[\frac{u_+}{\beta_N}\Q_{m_+(M)}(\beta_N)+e_+\right],
\]
which is \eqref{eq:central-upper}. The program has $B$ live roles and at most $D_0m_+(M)$ terminal roles. Lemma~\ref{lem:roles} gives exactly the sufficient widths in \eqref{eq:widths}. A checkpoint counts only the terminal width. No physical renewal law was used: the conditional hazard controls survival even if the hidden state after rejection differs arbitrarily from its former state.

If $\beta_N<1$, commit after the first attempt, and use an optimizer at cap one. The loss certificate with $q=1$ gives $u_+\Q_{m_+(M)}(1)+e_+$, and the whole-attempt cost is at most $\kappa_+$. This proves all remaining assertions.
\end{proof}

\section{Scope and further mathematical questions}
\label{sec:r73-scope}\label{sec:r74-scope}\label{sec:r75-scope}\label{sec:r76-scope}\label{sec:r77-scope}\label{sec:r78-scope}
The shared-preparation normal form applies to every finite group, every
finite preparation and outcome alphabet, and closed convex wordwise
specifications.  The input controller need not satisfy group relations.
The invariant polytope may extend beyond the target response span.
The exact low-contrast formula requires a common mean law; the displayed
nonnegative-rank counterexample explains why this hypothesis cannot
simply be removed.  These are shared-device results with one fixed
decoder, not a claim that independently optimized devices can always
share their memory at no cost.

The reset frontier now allows positive semidefinite failure effects
and classifies every live-label budget, with arbitrary ties.  Its normal
form uses the all-initial-state quantifier, a fixed row program, reported
resetting, a positive cost floor, and mandatory success stopping.  It
does not apply unchanged to a restricted initial prior, nonreset
back-action, extra observation ports or optional commitment.  Positive
exposure remains necessary for the separate simple transversal formula;
the general frontier correctly includes its semidefinite counterexample.
The old first-residual certificate remains a lower bound, and the exact
hierarchy exhibits a strict factor-two difference from the full cost gap.
Algebraic computability is not an efficient algorithm, and the inherited
covering reduction prevents a general polynomial-time memory-synthesis
claim unless the usual complexity obstruction is resolved.

The new approximation lower permits arbitrary classical terminal
decoder tables but requires delayed terminal test requests and no
uncounted access to the preparation history.  At fixed positive target
budget and the original algebraic phase it leaves a gap between
logarithmic-in-logarithm and logarithmic memory growth.  The sharper
matching logarithmic physical-dictionary law retains its Born-readout
and badly-approximable hypotheses.  Neither law is the original
general-$M$ squared-sign minimax-risk frontier.  The latter, general
nonrigid optimal feedback without the reset hypotheses,
and the intermediate phase-type timing frontier remain distinct
mathematical problems.

The inherited mean--contrast law is exact for its displayed growing
symmetric-group family, and the general finite-group error principle
is an attained finite optimization.  No closed formula for every group
or universal polynomial-time synthesis is inferred.  The return-polytope
construction concerns explicitly verified conditional branches, not
average mixing.  Its target-cost bound does not control the simulator's
unbounded expected duration.  The stationary-initialization theorem
concerns exact predictively stationary laws; an approximate truncation
does not automatically inherit its extra-label lower bound.

All original causal, predictive-state, information-compression, finite
resource, finite-order measurement, two-stage reset, rank-constrained
discrimination, communication, dual, spectral and geometric developments
retain their hypotheses.  Local geometry remains local.  Independent
A1, A2 and further Theta construction problems require their own
arguments.  Program bits, arithmetic precision, working storage and
sampling work are not silently identified with persistent label counts.
The paired supplement contains complete technical and independent
proofs; the separately compilable R58 companion retains its own scope.

\renewcommand{\bibliofont}{\footnotesize\setlength{\baselineskip}{10.8pt}}
\begin{thebibliography}{99}
\bibitem{Blackwell} D. Blackwell, \emph{Equivalent comparisons of experiments}, Ann. Math. Statist. \textbf{24} (1953), 265--272. DOI: 10.1214/aoms/1177729032.
\bibitem{CGM97} J. A. Cuesta-Albertos, A. Gordaliza and C. Matr\'an, \emph{Trimmed $k$-means: an attempt to robustify quantizers}, Ann. Statist. \textbf{25} (1997), 553--576. DOI: 10.1214/aos/1031833664.
\bibitem{GG99} L. A. Garc\'ia-Escudero and A. Gordaliza, \emph{Robustness properties of $k$ means and trimmed $k$ means}, J. Amer. Statist. Assoc. \textbf{94} (1999), 956--969. DOI: 10.1080/01621459.1999.10474200.
\bibitem{SC} C. R. Shalizi and J. P. Crutchfield, \emph{Computational mechanics: pattern and prediction, structure and simplicity}, J. Statist. Phys. \textbf{104} (2001), 817--879. DOI: 10.1023/A:1010388907793.
\bibitem{BC} N. Barnett and J. P. Crutchfield, \emph{Computational mechanics of input--output processes: structured transformations and the $\varepsilon$-transducer}, J. Statist. Phys. \textbf{161} (2015), 404--451; arXiv:1412.2690.
\bibitem{AMTW} A. Ambainis, M. Mosca, A. Tapp and R. de Wolf, \emph{Private quantum channels}, Proc. 41st Annual Symposium on Foundations of Computer Science (2000), 547--553. DOI: 10.1109/SFCS.2000.892142.
\bibitem{Buscemi} F. Buscemi, \emph{Comparison of quantum statistical models: equivalent conditions for sufficiency}, Comm. Math. Phys. \textbf{310} (2012), 625--647. DOI: 10.1007/s00220-012-1421-3.
\bibitem{BR14} N. Bertschinger and J. Rauh, \emph{The Blackwell relation defines no lattice}, Proc. IEEE International Symposium on Information Theory (2014), 2479--2483. DOI: 10.1109/ISIT.2014.6875280; arXiv:1401.3146.
\bibitem{Banerjee05} A. Banerjee, X. Guo and H. Wang, \emph{On the optimality of conditional expectation as a Bregman predictor}, IEEE Trans. Inform. Theory \textbf{51} (2005), 2664--2669. DOI: 10.1109/TIT.2005.850145.
\bibitem{Stengel96} B. von Stengel, \emph{Efficient computation of behavior strategies}, Games Econom. Behav. \textbf{14} (1996), 220--246. DOI: 10.1006/game.1996.0050.
\bibitem{Balder89} E. J. Balder, \emph{On compactness of the space of policies in stochastic dynamic programming}, Stochastic Process. Appl. \textbf{32} (1989), 141--150. DOI: 10.1016/0304-4149(89)90058-6.
\bibitem{Barry14} J. Barry, D. T. Barry and S. Aaronson, \emph{Quantum partially observable Markov decision processes}, Phys. Rev. A \textbf{90} (2014), 032311. DOI: 10.1103/PhysRevA.90.032311; arXiv:1406.2858.
\bibitem{Ying18} S. Ying and M. Ying, \emph{Reachability analysis of quantum Markov decision processes}, Inform. Comput. \textbf{263} (2018), 31--51. DOI: 10.1016/j.ic.2018.09.001; arXiv:1406.6146.
\bibitem{Saldi25} N. Saldi, S. Sanjari and S. Y\"uksel, \emph{Quantum Markov decision processes: general theory, approximations, and classes of policies}, SIAM J. Control Optim. \textbf{63} (2025), 1156--1183. DOI: 10.1137/24M1643864; arXiv:2402.14649.
\bibitem{Seidenberg54} A. Seidenberg, \emph{A new decision method for elementary algebra}, Ann. of Math. (2) \textbf{60} (1954), 365--374. DOI: 10.2307/1969640.
\bibitem{Chatterjee15} K. Chatterjee, M. Chmel\'ik, R. Gupta and A. Kanodia, \emph{Optimal cost almost-sure reachability in POMDPs}, Proc. AAAI Conference on Artificial Intelligence \textbf{29} (2015). DOI: 10.1609/aaai.v29i1.9683.
\bibitem{Andriushchenko23} R. Andriushchenko, A. Bork, M. \v{C}e\v{s}ka, S. Junges, J.-P. Katoen and M. Mac\'ak, \emph{Search and explore: symbiotic policy synthesis in POMDPs}, Computer Aided Verification (2023); arXiv:2305.14149v2.
\bibitem{Halmos69} P. R. Halmos, \emph{Two subspaces}, Trans. Amer. Math. Soc. \textbf{144} (1969), 381--389. DOI: 10.1090/S0002-9947-1969-0251519-5.
\bibitem{Kessler21} D. A. Kessler, E. Barkai and K. Ziegler, \emph{First-detection time of a quantum state under random probing}, Phys. Rev. A \textbf{103} (2021), 022222. DOI: 10.1103/PhysRevA.103.022222; arXiv:2012.01763.
\bibitem{DelVecchio26} G. Del Vecchio Del Vecchio and S. N. Majumdar, \emph{Optimal detection of quantum states via projective measurements}, J. Phys. A: Math. Theor. \textbf{59} (2026), 035001. DOI: 10.1088/1751-8121/ae34fe; arXiv:2509.08556.
\bibitem{Companion} Q. Qi, \emph{General Theta Foundations I: Acquired Geometry and Causal Resource Transfer}, complete technical companion on absolute calibration, exact stopping resources and nonintervening sensing, 2026. The complete retained R58 native source and its independently built companion accompany this manuscript; their identities are fixed in the source-provenance and preservation records.
\bibitem{Tzeng92} W.-G. Tzeng, \emph{A polynomial-time algorithm for the equivalence of probabilistic automata}, SIAM J. Comput. \textbf{21} (1992), 216--227. DOI: 10.1137/0221017.
\bibitem{LiYuYing12} Y. Li, N. Yu and M. Ying, \emph{Termination of nondeterministic quantum programs}, arXiv:1201.0891 (2012).
\bibitem{Fu24} J. Fu, H. Jiang, M. Xu, Y. Deng and Z.-B. Li, \emph{Algorithmic analysis of termination problems for nondeterministic quantum programs}, arXiv:2402.15827 (2024).

\bibitem{LiuVerdu18} J. Liu and S. Verd\'u, \emph{Rejection sampling and noncausal sampling under moment constraints}, Proc. IEEE International Symposium on Information Theory (2018), 1565--1569. DOI: 10.1109/ISIT.2018.8437857.
\bibitem{GocFlamich24} D. Goc and G. Flamich, \emph{On channel simulation with causal rejection samplers}, Proc. IEEE International Symposium on Information Theory (2024). DOI: 10.1109/ISIT57864.2024.10619339; examined text arXiv:2401.16579v2 (3 May 2024).
\bibitem{LiLiu26} Z. Li and T. Liu, \emph{Ranking experiments under sequential sampling}, arXiv:2608.19897v2 (28 August 2026). Version 1 appeared under the title \emph{The order of binary experiments under endogenous stopping}, with Z. Li as sole author; the version-2 statement is the one compared here.
\bibitem{Ying14pre} S. Ying and M. Ying, \emph{Reachability analysis of quantum Markov decision processes}, arXiv:1406.6146 (2014), Theorem 3.6; published journal version listed above as \cite{Ying18}.

\bibitem{KaraYuksel20} A. D. Kara and S. Y\"uksel, \emph{Robustness to incorrect system models in stochastic control}, SIAM J. Control Optim. \textbf{58} (2020), 1144--1182. DOI: 10.1137/18M1208058; examined text arXiv:1803.06046v7.
\bibitem{Demirci25} Y. E. Demirci, A. D. Kara and S. Y\"uksel, \emph{Sensitivity of filter kernels and robustness bounds to transition and measurement kernel perturbations in partially observable stochastic control}, arXiv:2508.10658v3 (2025).
\bibitem{Watrous18} J. Watrous, \emph{The Theory of Quantum Information}, Cambridge University Press, 2018, Chapter 3, Corollary 3.40.
\bibitem{KRN65} K. Kuratowski and C. Ryll-Nardzewski, \emph{A general theorem on selectors}, Bull. Acad. Polon. Sci. S\'er. Sci. Math. Astronom. Phys. \textbf{13} (1965), 397--403.

\bibitem{Charnes63} A. Charnes, W. W. Cooper and K. O. Kortanek, \emph{Duality in semi-infinite programs and some works of Haar and Carath\'eodory}, Management Sci. \textbf{9} (1963), 209--228. DOI: 10.1287/mnsc.9.2.209.
\bibitem{Shapiro09} A. Shapiro, \emph{Semi-infinite programming, duality, discretization and optimality conditions}, Optimization \textbf{58} (2009), 133--161. DOI: 10.1080/02331930902730070; author preprint, Optimization Online (2008).
\bibitem{Dufour25} F. Dufour and T. Prieto-Rumeau, \emph{Absorbing Markov decision processes and their occupation measures}, SIAM J. Control Optim. \textbf{63} (2025), 676--698. DOI: 10.1137/24M1661959. Related author preprint examined: \emph{A note on weak compactness of occupation measures for an absorbing Markov decision process}, arXiv:2402.10672v2 (2024).
\bibitem{Franc23} V. Franc, D. Prusa and V. Voracek, \emph{Optimal strategies for reject option classifiers}, J. Mach. Learn. Res. \textbf{24} (2023), no. 11, 1--49.
\bibitem{Jencova16} A. Jen\v{c}ov\'a, \emph{Comparison of quantum channels and statistical experiments}, arXiv:1512.07016v2 (2016). The separately indexed shorter version arXiv:1601.06370 was withdrawn in favor of this extended text.
\bibitem{LiElGamal18} C. T. Li and A. El Gamal, \emph{Strong functional representation lemma and applications to coding theorems}, IEEE Trans. Inform. Theory \textbf{64} (2018), 6967--6978; arXiv:1701.02827v4.
\bibitem{AldousShepp87} D. Aldous and L. Shepp, \emph{The least variable phase type distribution is Erlang}, Comm. Statist. Stochastic Models \textbf{3} (1987), 467--473. DOI: 10.1080/15326348708807067.
\bibitem{Sion58} M. Sion, \emph{On general minimax theorems}, Pacific J. Math. \textbf{8} (1958), 171--176. DOI: 10.2140/pjm.1958.8.171.
\bibitem{PiunovskiyZhang24} A. Piunovskiy and Y. Zhang, \emph{Extreme occupation measures in Markov decision processes with an absorbing state}, SIAM J. Control Optim. \textbf{62} (2024), 65--90. DOI: 10.1137/23M1572398. Author text examined: \emph{Extreme occupation measures in Markov decision processes with a cemetery}, arXiv:2307.03158v1.
\bibitem{DufourMixtures} F. Dufour and T. Prieto-Rumeau, \emph{Absorbing Markov decision processes: geometric properties and sufficiency of finite mixtures of deterministic policies}, arXiv:2512.17511; examined HTML at version-1 URL, Theorems 4.3--4.4 and Proposition 4.5.
\bibitem{Mereghetti24} C. Mereghetti, B. Palano and P. Raucci, \emph{Unary quantum finite state automata with control language}, Appl. Sci. \textbf{14} (2024), no. 4, 1490. DOI: 10.3390/app14041490.
\bibitem{RabinScott59} M. O. Rabin and D. Scott, \emph{Finite automata and their decision problems}, IBM J. Res. Develop. \textbf{3} (1959), 114--125. DOI: 10.1147/rd.32.0114.
\bibitem{Jaeger00} H. Jaeger, \emph{Observable operator models for discrete stochastic time series}, Neural Comput. \textbf{12} (2000), 1371--1398. DOI: 10.1162/089976600300015411.
\bibitem{ThonJaeger15} M. Thon and H. Jaeger, \emph{Links between multiplicity automata, observable operator models and predictive state representations---a unified learning framework}, J. Mach. Learn. Res. \textbf{16} (2015), 103--147.
\bibitem{MonrasWinter14} A. Monras and A. Winter, \emph{Quantum learning of classical stochastic processes: the completely-positive realization problem}, LIPIcs \textbf{27}, TQC 2014 (2014), 99--109. DOI: 10.4230/LIPIcs.TQC.2014.99.

\bibitem{Heller65} A. Heller, \emph{On stochastic processes derived from Markov chains}, Ann. Math. Statist. \textbf{36} (1965), 1286--1291. DOI: 10.1214/aoms/1177700000.
\bibitem{CarlylePaz71} J. W. Carlyle and A. Paz, \emph{Realizations by stochastic finite automata}, J. Comput. System Sci. \textbf{5} (1971), 26--40. DOI: 10.1016/S0022-0000(71)80005-3.
\bibitem{Fanizza24} M. Fanizza, J. Lumbreras and A. Winter, \emph{Quantum theory in finite dimension cannot explain every general process with finite memory}, Comm. Math. Phys. \textbf{405} (2024), article 50. DOI: 10.1007/s00220-023-04913-4; arXiv:2209.11225. See Section 2.3 for the FRDN phase examples and Theorem 2 and Section 4.2 for the noise-dependent worst-phase classical-memory lower bound.
\bibitem{DmitrievDynkin46} N. Dmitriev and E. Dynkin, \emph{On characteristic roots of stochastic matrices}, Izv. Akad. Nauk SSSR Ser. Mat. \textbf{10} (1946), no. 2, 167--184. Math-Net.Ru: im3595 (Russian).
\bibitem{JohnsonPaparella17} C. R. Johnson and P. Paparella, \emph{A matricial view of the Karpelevi\v c theorem}, Linear Algebra Appl. \textbf{520} (2017), 1--15. DOI: 10.1016/j.laa.2017.01.009; arXiv:1611.06970v2.
\bibitem{KirklandSmigoc22} S. Kirkland and H. \v Smigoc, \emph{Stochastic matrices realising the boundary of the Karpelevi\v c region}, Linear Algebra Appl. \textbf{635} (2022), 116--138. DOI: 10.1016/j.laa.2021.11.016; arXiv:2110.01040v1.
\bibitem{BenvenutiFarina04} L. Benvenuti and L. Farina, \emph{A tutorial on the positive realization problem}, IEEE Trans. Automat. Control \textbf{49} (2004), 651--664. DOI: 10.1109/TAC.2004.826715.
\bibitem{BenvenutiFarina99} L. Benvenuti and L. Farina, \emph{An example of how positivity may force realizations of ``large'' dimension}, Systems Control Lett. \textbf{36} (1999), 261--266. DOI: 10.1016/S0167-6911(98)00098-X.
\bibitem{OCinneide90} C. A. O'Cinneide, \emph{Characterization of phase-type distributions}, Comm. Statist. Stochastic Models \textbf{6} (1990), 1--57. DOI: 10.1080/15326349908807134.
\bibitem{Maier91} R. S. Maier, \emph{The algebraic construction of phase-type distributions}, Comm. Statist. Stochastic Models \textbf{7} (1991), 573--602. DOI: 10.1080/15326349108807207.
\bibitem{HoneVarona19} A. N. W. Hone and J. L. Varona, \emph{Continued fractions and irrationality exponents for modified Engel and Pierce series}, Monatsh. Math. \textbf{190} (2019), 501--516. DOI: 10.1007/s00605-018-1244-1. Author version arXiv:1810.10127v1, equation (3.3).
\bibitem{Nussbaum98} R. D. Nussbaum and M. Scheutzow, \emph{Admissible arrays and a nonlinear generalization of Perron--Frobenius theory}, J. London Math. Soc. (2) \textbf{58} (1998), 526--544. DOI: 10.1112/S0024610798006796.
\bibitem{Brown64} D.~R. Brown, On clans of non-negative matrices,
\emph{Proc. Amer. Math. Soc.} \textbf{15} (1964), 671--674,
doi:10.1090/S0002-9939-1964-0165504-6.
\bibitem{Schwarz64} \v{S}. Schwarz, On the structure of the semigroup of
stochastic matrices, \emph{Publ. Math. Inst. Hungar. Acad. Sci.}
\textbf{9} (1964), 297--311, especially Theorems~7--8.
\bibitem{Flor69} P. Flor, On groups of non-negative matrices,
\emph{Compositio Math.} \textbf{21} (1969), 376--382,
\url{https://www.numdam.org/article/CM_1969__21_4_376_0.pdf}.
\bibitem{Serre07} J.-P. Serre, Bounds for the orders of the finite subgroups
of $G(k)$, in \emph{Group Representation Theory}, M. Geck, D. Testerman
and J. Th\'evenaz (eds.), EPFL Press, Lausanne, 2007, 405--450;
author lecture text, April 2006, archived as arXiv:1011.0346.
\bibitem{BabaiBealsRockmore93} L. Babai, R. Beals and D. Rockmore,
Deciding finiteness of matrix groups in deterministic polynomial time,
\emph{Proceedings of ISSAC 1993}, 117--126,
doi:10.1145/164081.164104.
\bibitem{KieferWachter14} S. Kiefer and B. Wachter, Stability and complexity
of minimising probabilistic automata, in \emph{ICALP 2014, Part II},
Lecture Notes in Computer Science \textbf{8573}, 268--279,
doi:10.1007/978-3-662-43951-7\_23; arXiv:1404.6673, Section~4.2.
\bibitem{Schwarz62} \v{S}. Schwarz, Homomorphisms of a completely simple
semigroup onto a group, \emph{Matematicko-fyzik\'alny \v{c}asopis}
\textbf{12} (1962), no.~4, 293--300;
\url{https://www.dml.cz/handle/10338.dmlcz/126333}.

\bibitem{GPT13} J. Gouveia, P. A. Parrilo and R. R. Thomas,
\emph{Lifts of convex sets and cone factorizations}, Math. Oper. Res.
\textbf{38} (2013), 248--264. DOI: 10.1287/moor.1120.0575;
arXiv:1111.3164, especially Theorem 2.12 on equivariant factorizations.
\bibitem{KPT10} V. Kaibel, K. Pashkovich and D. O. Theis,
\emph{Symmetry matters for sizes of extended formulations},
SIAM J. Discrete Math. \textbf{26} (2012), 1361--1382.
DOI: 10.1137/110839813; arXiv:0911.3712v3.
An extended abstract appeared in IPCO 2010, LNCS 6080, 135--148.
\bibitem{BocciCarliniRapallo11} C. Bocci, E. Carlini and F. Rapallo,
\emph{Perturbation of matrices and nonnegative rank with a view toward
statistical models}, SIAM J. Matrix Anal. Appl. \textbf{32} (2011),
1500--1512. DOI: 10.1137/110825455; arXiv:1010.3566.
\bibitem{GillisGlineur12} N. Gillis and F. Glineur,
\emph{On the geometric interpretation of the nonnegative rank},
Linear Algebra Appl. \textbf{437} (2012), 2685--2712.
DOI: 10.1016/j.laa.2012.06.038; arXiv:1009.0880.
\bibitem{Okay26} C. Okay, \emph{Polyhedral classical simulators for quantum
computation}, Res. Math. Sci. \textbf{13} (2026), article 55.
DOI: 10.1007/s40687-026-00626-x. In particular Theorem 16, Lemma 17
and Corollary 19 treat invariant-polytope updates and adaptive composition.
\bibitem{AubrunSzarek17} G. Aubrun and S. J. Szarek,
\emph{Dvoretzky's theorem and the complexity of entanglement detection},
Discrete Anal. (2017), paper no. 1, 20 pp.; arXiv:1510.00578,
especially Section 6 on polyhedral approximations of density matrices.
\bibitem{Fawzi21} H. Fawzi,
\emph{On polyhedral approximations of the positive semidefinite cone},
Math. Oper. Res. \textbf{46} (2021), 1479--1489.
DOI: 10.1287/moor.2020.1077; arXiv:1811.09649.
\bibitem{HolroydLiggett16} A. E. Holroyd and T. M. Liggett,
\emph{Finitely dependent coloring}, Forum Math. Pi \textbf{4} (2016),
e9, 43 pp. DOI: 10.1017/fmp.2016.7; arXiv:1403.2448.
See Section 5 and the finite hidden-Markov obstruction attributed there
to Schramm.
\bibitem{Serre77} J.-P. Serre, \emph{Linear Representations of Finite Groups},
Graduate Texts in Mathematics \textbf{42}, Springer-Verlag, New York, 1977.
DOI: 10.1007/978-1-4684-9458-7.

\bibitem{MonrasWinter16} A. Monras and A. Winter,
\emph{Quantum learning of classical stochastic processes: the
completely-positive realization problem}, J. Math. Phys. \textbf{57}
(2016), 015219. DOI: 10.1063/1.4936935; arXiv:1412.3634.

\bibitem{VuThanhGillisLecron23} O. Vu Thanh, N. Gillis and F. Lecron,
\emph{Bounded simplex-structured matrix factorization: algorithms,
identifiability and applications}, IEEE Trans. Signal Process.
\textbf{71} (2023), 2434--2447.
DOI: 10.1109/TSP.2023.3289704; arXiv:2209.12638v2,
especially Section IV-C and Theorem 3.
\bibitem{Gillis12} N. Gillis,
\emph{Sparse and unique nonnegative matrix factorization through data
preprocessing}, J. Mach. Learn. Res. \textbf{13} (2012), 3349--3386,
especially Section 4.3, Lemma 27, Example 2 and Theorem 29.
\url{https://jmlr.org/papers/v13/gillis12a.html}.

\bibitem{Rabin63} M. O. Rabin, \emph{Probabilistic automata},
Information and Control \textbf{6} (1963), 230--245.
DOI: 10.1016/S0019-9958(63)90290-0, especially Theorem 3 and Section VIII.

\bibitem{Li80}
S.-Y. R. Li, \emph{A martingale approach to the study of occurrence of
sequence patterns in repeated experiments}, Ann. Probab. \textbf{8}
(1980), no. 6, 1171--1176, doi:10.1214/aop/1176994578.


\bibitem{BertsekasTsitsiklis91} D. P. Bertsekas and J. N. Tsitsiklis,
\emph{An analysis of stochastic shortest path problems},
Math. Oper. Res. \textbf{16} (1991), 580--595.
DOI: \texttt{10.1287/moor.16.3.580}.
\bibitem{Meuleau99} N. Meuleau, K.-E. Kim, L. P. Kaelbling and A. R. Cassandra,
\emph{Solving POMDPs by searching the space of finite policies},
Proceedings of UAI (1999), 417--426; author's deposit arXiv:1301.6720.
\bibitem{PoupartBoutilier03} P. Poupart and C. Boutilier,
\emph{Bounded finite state controllers},
Advances in Neural Information Processing Systems \textbf{16} (2003).
\bibitem{Karp72} R. M. Karp,
\emph{Reducibility among combinatorial problems},
in R. E. Miller and J. W. Thatcher (eds.), Complexity of Computer
Computations, Plenum Press (1972), 85--103.
DOI: \texttt{10.1007/978-1-4684-2001-2\_9}.
\bibitem{VandaeleGillisGlineur17} A. Vandaele, N. Gillis and F. Glineur,
\emph{On the linear extension complexity of regular $n$-gons},
Linear Algebra Appl. \textbf{521} (2017), 217--239.
DOI: 10.1016/j.laa.2016.12.023; arXiv:1505.08031v2,
especially Section 3.1 on support antichains and rectangle covering.

% Verified primary-source metadata, 23 October 2026.
% Existing BertsekasTsitsiklis91, Meuleau99 and PoupartBoutilier03 entries
% should remain; these are proposed additions, not a separate bibliography.
\bibitem{AmatoBernsteinZilberstein10}
C. Amato, D. S. Bernstein and S. Zilberstein,
Optimizing fixed-size stochastic controllers for POMDPs and decentralized
POMDPs, \emph{Autonomous Agents and Multi-Agent Systems} \textbf{21}
(2010), 293--320. \url{https://doi.org/10.1007/s10458-009-9103-z}.


\end{thebibliography}

\end{document}
